\documentclass[11pt,oneside]{article}

\usepackage[T1]{fontenc}
\usepackage[utf8]{inputenc}
\usepackage{lmodern}
\usepackage{microtype}
\usepackage[a4paper,margin=25mm,headheight=15pt]{geometry}
\usepackage{amsmath,amssymb,amsthm,mathtools}
\usepackage{booktabs,longtable,tabularx,array,multirow}
\usepackage{enumitem}
\usepackage[dvipsnames,table]{xcolor}
\usepackage{listings}
\usepackage{fancyhdr}
\usepackage{hyperref}
\usepackage[nameinlink,noabbrev]{cleveref}
\usepackage{bookmark}
\usepackage{url}

\hypersetup{
  colorlinks=true,
  linkcolor=MidnightBlue,
  citecolor=OliveGreen,
  urlcolor=RoyalBlue,
  pdftitle={Further Progress on the Alaoglu--Erdos Two-Base Exponent Conjecture},
  pdfauthor={Research report prepared for Vladimir Reshetnikov},
  pdfsubject={Exact conditional structure, new equivalent formulations, and certified finite exclusion}
}
\urlstyle{same}

\pagestyle{fancy}
\fancyhf{}
\fancyhead[L]{Alaoglu--Erd\H{o}s: further progress}
\fancyhead[R]{\thepage}
\renewcommand{\headrulewidth}{0.4pt}

\setlength{\parindent}{0pt}
\setlength{\parskip}{5.5pt plus 1.5pt minus 1pt}
\setlength{\emergencystretch}{3em}
\setlist[itemize]{topsep=3pt,itemsep=2pt,parsep=0pt,leftmargin=1.6em}
\setlist[enumerate]{topsep=3pt,itemsep=2pt,parsep=0pt,leftmargin=1.9em}
\allowdisplaybreaks

\newtheorem{theorem}{Theorem}[section]
\newtheorem{proposition}[theorem]{Proposition}
\newtheorem{lemma}[theorem]{Lemma}
\newtheorem{corollary}[theorem]{Corollary}
\newtheorem{conjecture}[theorem]{Conjecture}
\theoremstyle{definition}
\newtheorem{definition}[theorem]{Definition}
\newtheorem{remark}[theorem]{Remark}
\newtheorem{question}[theorem]{Research question}

\newcommand{\N}{\mathbb N}
\newcommand{\Nzero}{\mathbb N_0}
\newcommand{\Z}{\mathbb Z}
\newcommand{\Q}{\mathbb Q}
\newcommand{\R}{\mathbb R}
\newcommand{\C}{\mathbb C}
\newcommand{\Qbar}{\overline{\mathbb Q}}
\newcommand{\Sol}{\mathcal S}
\newcommand{\Cand}{\mathcal C}
\newcommand{\Smooth}{\mathcal H_{2,3}}
\newcommand{\Sat}{\operatorname{Sat}}
\newcommand{\supp}{\operatorname{supp}}
\newcommand{\clog}{\operatorname{clog}}
\newcommand{\vpp}[2]{v_{#1}\!\left(#2\right)}
\newcommand{\tagLean}{\textcolor{ForestGreen}{\textsf{[kernel-verified Lean]}}}
\newcommand{\tagPaper}{\textcolor{MidnightBlue}{\textsf{[paper proof in this report]}}}
\newcommand{\tagClassical}{\textcolor{Purple}{\textsf{[classical theorem used]}}}
\newcommand{\tagComp}{\textcolor{BurntOrange}{\textsf{[directed-rounding computation]}}}
\newcommand{\tagProposal}{\textcolor{BrickRed}{\textsf{[research direction]}}}
\newcommand{\lean}[1]{\texttt{\detokenize{#1}}}
\newcommand{\commit}[1]{\texttt{#1}}

\newcolumntype{Y}{>{\raggedright\arraybackslash}X}
\newcolumntype{Z}{>{\centering\arraybackslash}p{24mm}}

\lstset{
  basicstyle=\ttfamily\scriptsize,
  keywordstyle=\bfseries,
  commentstyle=\itshape,
  columns=fullflexible,
  keepspaces=true,
  showstringspaces=false,
  breaklines=true,
  breakatwhitespace=false,
  frame=single,
  rulecolor=\color{black!25},
  xleftmargin=2mm,
  xrightmargin=2mm,
  aboveskip=6pt,
  belowskip=6pt,
  literate={·}{{$\cdot$}}1
}

\title{\bfseries Further Progress on the Alaoglu--Erd\H{o}s\\
Two-Base Exponent Conjecture\\[5pt]
\large Multiplicative rigidity, transcendental product exponents,\\
primitive-power statistics, ambient root saturation,\\
and a certified search through $2^{26}$}
\author{Research report prepared for Vladimir Reshetnikov}
\date{20 August 2026}

\begin{document}
\maketitle

\begin{abstract}
Let
\[
  \Sol=\{x\in\R:2^x\in\Z,\ 3^x\in\Z\}.
\]
The Alaoglu--Erd\H{o}s conjecture asserts that \(\Sol=\Nzero\).  This report starts from
the Lean development and the combined report in the \texttt{ProveIt} repository at commit
\commit{5d3982ceb12644fdd3499569d9db5f62f9a31dc6}.  No unconditional proof of the
conjecture is obtained.  Its final obstruction remains a special case of the unresolved
four-exponentials problem.  The investigation nevertheless produces several new exact
consequences and equivalent formulations.

Conditional on a counterexample, the repository proves that there is a transcendental least
noninteger solution \(\beta\) for which every solution has unique coordinates
\[
  \Sol=\{n+k\beta:n,k\in\Nzero\}.
\]
We prove that this monoid is rigid under every rational function with algebraic
coefficients: if \(R\in\Qbar(T)\) is regular at \(\beta\), then
\(R(\beta)\in\Sol\) exactly when \(R(T)=n+kT\) for some \(n,k\in\Nzero\).
In particular, for arbitrary \(x,y\in\Sol\),
\[
  xy\in\Sol\quad\Longleftrightarrow\quad x\in\Nzero\ \text{or}\ y\in\Nzero.
\]
Thus the conjecture is equivalent to multiplicative closure of \(\Sol\), and even to
closure under squaring alone.

A second argument combines the repository's three-smooth-base theorem with the classical
six-exponentials theorem.  If \(x,y\in\Sol\) are both nonintegral, then
\(6^{xy}\) is transcendental; each of \(2^{xy}\) and \(3^{xy}\) is either an
integer or transcendental, and at least one is transcendental.  Hence failure of
multiplicative closure is not a one-unit integrality miss: it is witnessed by genuine
transcendence.

The exact coordinates also classify every common perfect-power and affine-descent pattern.
If \(x=n+k\beta\), the two outputs are simultaneously perfect \(r\)-th powers exactly
when \(r\mid n\) and \(r\mid k\); their maximal common power degree is
\(\gcd(n,k)\).  Consequently the proportion of simultaneously power-primitive output
pairs of height below \(T\) tends to \(1/\zeta(2)=6/\pi^2\), and the exact degree
\(g\) occurs with limiting probability \(6/(\pi^2g^2)\).  We further identify the
candidate divisibility lattice, prove unconditional closure under gcd and lcm, enumerate
candidate divisors, and compute the ambient root saturation in terms of the primitive odd
core and the rational-power index.

Finally, a five-thread MPFR~4.2.2 verifier using 192-bit directed rounding checks every
non-power of two below \(2^{26}\), a total of \(67{,}108{,}837\) inputs.  No candidate
occurs.  At the explicitly stated software trust level, every hypothetical noninteger
solution therefore satisfies \(x>26\).  A variable-order optimization of the formalized
finite-difference obstruction explains quantitatively why that method, by itself, still
falls short.  The report ends with a prioritized formalization and research program.
\end{abstract}

\tableofcontents
\newpage

\section{Executive summary}\label{sec:executive}

\subsection{Bottom line}

The conjecture remains unproved.  This conclusion is not a failure of effort but a precise
mathematical diagnosis.  The repository already reduces any counterexample to a canonical
primitive radical and a free rank-two solution monoid.  Every elementary counting,
progression, descent, and equidistribution constraint currently known is internally
consistent with that monoid.  To finish the proof one needs a new input at the
four-exponentials boundary, or an argument that uses the special \emph{integrality and
valuation structure} more strongly than general transcendence theory presently does.

The strongest additions of this report are:
\begin{enumerate}
\item an exact rational-function intersection theorem for the canonical generator;
\item unconditional multiplication, quotient, and power criteria inside \(\Sol\);
\item algebraicity and transcendence reformulations based on six exponentials;
\item a complete common-perfect-power and matched-affine-decomposition classification;
\item primitive-output asymptotics and exact limiting degree statistics;
\item a candidate gcd/lcm lattice and exact ambient root saturation;
\item the optimized variable-order scale \(n/(\theta\log n)\) for the higher-difference
method, together with a sharp quantitative explanation of its gap;
\item a directed-rounding finite exclusion through \(2^{26}\).
\end{enumerate}

\subsection{Trust labels}

The report keeps four proof levels separate.
\begin{center}
\begin{tabularx}{0.95\textwidth}{@{}Z Y@{}}
\toprule
\tagLean & Accepted by the Lean kernel in the inspected repository snapshot.\\
\tagPaper & A conventional proof is supplied here; it has not yet been checked by Lean.\\
\tagClassical & The proof invokes a named published theorem, principally six exponentials.\\
\tagComp & A finite statement checked by directed-rounding software; this is not a kernel proof.\\
\bottomrule
\end{tabularx}
\end{center}

The finite search is useful evidence and a stronger explicit bound, but it does not change
the logical status of the infinite conjecture.  The kernel-verified bound remains
\(x\ge12\) for a noninteger solution; the new \(x>26\) bound carries the disclosed
software trust boundary.

\section{The problem and the audited repository snapshot}\label{sec:problem}

\subsection{Original and candidate forms}

\begin{conjecture}[Alaoglu--Erd\H{o}s, two-base integer form]\label{conj:AE}
For every real number \(x\),
\[
  2^x\in\Z\quad\text{and}\quad 3^x\in\Z
  \quad\Longrightarrow\quad x\in\Z.
\]
\end{conjecture}

The hypotheses force \(x\ge0\), so the intended conclusion is \(x\in\Nzero\).  Put
\[
  \theta:=\log_2 3=\frac{\log3}{\log2}
   =1.58496250072115618145\ldots
\]
and define the natural-candidate set
\[
  \Cand:=\{m\in\N_{>0}:m^\theta\in\N\}.
\]
The maps \(x\mapsto2^x\) and \(m\mapsto\log_2m\) give inverse bijections between
\(\Sol\) and \(\Cand\), because
\[
  (2^x)^\theta=3^x.
\]
Integral solutions correspond exactly to powers of two.  Thus
\[
  \text{Conjecture~\ref{conj:AE}}
  \quad\Longleftrightarrow\quad
  \Cand=\{2^n:n\in\Nzero\}.
\]

A positive natural number will be called \emph{three-smooth} when all its prime factors
belong to \(\{2,3\}\), equivalently when it has the form \(2^u3^v\).  The set of such
numbers is denoted \(\Smooth\).

\subsection{Repository snapshot}

The inspected branch is \texttt{main} at commit
\[
  \commit{5d3982ceb12644fdd3499569d9db5f62f9a31dc6},
\]
whose commit time is 21 August 2026 at 02:43:24 UTC, still 20 August in Pacific time.  The
starting document is
\begin{center}
\small
\texttt{Analysis/ExponentialIdentities/Article/}\\
\texttt{alaoglu-erdos-combined-report/}\\
\texttt{alaoglu-erdos-combined-report.tex}.
\end{center}
Its accompanying formal development lies under
\begin{center}
\small
\texttt{Analysis/ExponentialIdentities/Lean/}\\
\texttt{ExponentialIdentities/}.
\end{center}

The combined report is already a substantial synthesis: it records the three-base theorem,
Gelfond--Schneider consequences, exact finite certificates through \(4096\), rank-two and
common-odd-core structure, the canonical primitive generator, exact counting and Fourier
block averages, output normalization, localized-radical equivalences, exact radical degree,
and arbitrary-order finite-difference obstructions.  The present report treats those facts
as the established frontier and asks what further consequences can be forced from them.

\subsection{Why this is a four-exponentials problem}

If \(x\) were a noninteger solution, then \(1,x\) would be linearly independent over
\(\Q\), and \(\log2,\log3\) would also be linearly independent over \(\Q\).  Yet all
four exponentials
\[
  e^{\log2}=2,\qquad e^{\log3}=3,\qquad
  e^{x\log2}=2^x,\qquad e^{x\log3}=3^x
\]
would be algebraic.  The four-exponentials conjecture would rule this out immediately.
The six-exponentials theorem succeeds once a third independent logarithm or exponent is
available, but the \(2\times2\) case remains open
\cite{Ramachandra1968II,Roy1992,Waldschmidt2000,TheoremDB2026}.

This observation is a methodological firewall.  A proposed argument that uses only the
algebraicity of the four displayed values is, in effect, attempting a new proof of a
four-exponentials case.  A more plausible route must exploit features absent from the
general conjecture: positivity, integrality, prime valuations, exact denominator support,
or the canonical least-output normalization.

\section{The formalized frontier used in this report}\label{sec:audit}

This section states the repository results that enter the new arguments.  It is deliberately
selective: the combined report remains the detailed inventory, while the goal here is to
make every dependency of the new theorems visible.

\subsection{Elementary restrictions and transcendence}

The development proves the following.
\begin{itemize}
\item \tagLean If \(2^x\in\Z\), then \(x\ge0\).
\item \tagLean If \(x\in\Q\) and \(2^x\in\Z\), then \(x\in\Z\).  Hence every
noninteger solution is irrational.
\item \tagLean If \(2^x\in\Z\), then \(x\) is either an integer or transcendental over
\(\Q\).  In particular every noninteger two-base solution is transcendental.  The
headline theorem is
\lean{transcendental_of_not_integer_of_two_three_rpow_integer} in
\lean{Transcendence.lean}.
\item \tagLean The constant \(\theta=\log_2 3\) is transcendental.  Indeed
\(2^\theta=3\), while \(1<\theta<2\), so the same Gelfond--Schneider dichotomy applies.
\end{itemize}

The full Gelfond--Schneider theorem is formalized in the repository, following the Mathlib
formalization of Karatarakis.  The new arguments below use only the two consequences just
listed, not the internal auxiliary-function proof.

\subsection{Kernel-verified finite region}

\tagLean The finite-region proof is distributed across three modules:
\begin{center}
\small
\lean{FiniteCheck.lean},\qquad
\lean{FiniteCheck256.lean},\qquad
\lean{FiniteCheck4096.lean}.
\end{center}
Together they establish
\[
  m\in\Cand,\quad m<4096=2^{12}
  \quad\Longrightarrow\quad m\text{ is a power of }2.
\]
Equivalently,
\[
  x\in\Sol\setminus\Nzero\quad\Longrightarrow\quad x\ge12.
\]
The \(4096\) theorem is replayed by the kernel from a table of exact integer comparisons
using rational enclosures for \(\theta\).  It contains no \lean{sorry} and does not use
\lean{native_decide}.  This is a stronger trust model than an interval scan.

The combined report also described an independent software-level scan through \(2^{20}\).
\Cref{sec:computation} replaces it by a directed-rounding MPFR verifier and extends the
range sixty-fourfold, through \(2^{26}\).

\subsection{Rank two, the primitive odd core, and the exact solution monoid}

The three-base determinant theorem implies that prime-valuation vectors of candidates have
rational rank at most two.  Since \(2\) is already a candidate, every odd part lies on one
rational ray.  The repository normalizes that ray to a unique primitive odd core.

The decisive theorem for the present report is
\begin{center}
\small
\lean{exists_canonical_primitiveGenerator_of_not_alaogluErdosConjecture}
\end{center}
from \lean{PrimitiveGenerator.lean}.

\begin{theorem}[Canonical primitive generator, repository theorem]\label{thm:exact-monoid}
\tagLean Assume that the conjecture fails.  Then there exist natural numbers
\(w,d,a,A\) and a real number \(\beta\) such that:
\begin{enumerate}[label=\textup{(\roman*)}]
\item \(w>1\) is odd and its positive odd-prime valuations have gcd one;
\item \(d>0\), \(A>0\), and \(a=0\) or \(3\nmid A\);
\item with \(E=3^{\log_2w}=w^\theta\),
\[
  E^d=\frac{A}{3^a},
  \qquad
  E^j\in\Q\Longleftrightarrow d\mid j\quad(j\in\Z);
\]
\item
\[
  \beta=a+d\log_2w,\qquad
  M:=2^\beta=2^aw^d\in\N,\qquad
  3^\beta=A\in\N;
\]
\item \(\beta\) is the unique least noninteger solution and is irrational;
\item every solution has a unique coordinate pair
\[
  \boxed{\quad x=n+k\beta,\qquad n,k\in\Nzero;\quad}
\]
\item every candidate has a unique coordinate pair
\[
  \boxed{\quad m=2^nM^k,
    \qquad n,k\in\Nzero.\quad}
\]
\end{enumerate}
\end{theorem}

Combining the theorem with the formalized Gelfond--Schneider consequence gives:

\begin{corollary}\label{cor:beta-trans}
\tagLean Under failure, \(\beta\) is transcendental over \(\Q\), and hence also over
\(\Qbar\).
\end{corollary}

The last implication uses transitivity of algebraicity: a number algebraic over the
algebraic closure \(\Qbar\) would be algebraic over \(\Q\).

The theorem says much more than ``rank at most two.''  It identifies the entire solution
set with the free additive monoid on \(1,\beta\), the candidate set with the free
multiplicative monoid on \(2,M\), and the output-pair monoid with the free monoid on
\((2,3)\) and \((M,A)\).  There are no omitted lattice points or Frobenius gaps.

\subsection{Least-output and simultaneous normalization}

\tagLean The least solution is affinely primitive: if
\[
  M=2^rU^q,\qquad A=3^rV^q
\]
with \(q>0\), then \(r=0\) and \(q=1\).  In particular \(M,A\) are not both perfect
\(q\)-th powers for any \(q\ge2\), and not both divisible by their distinguished bases.

The stronger simultaneous normalization in \lean{OutputNormalization.lean} writes
\[
  M=2^aw^d,
  \qquad
  A=3^bv^e,
\]
where \(w\) is odd and power-primitive, \(v\) is not divisible by \(3\) and is
power-primitive, \(d,e>0\), and
\[
  a=0\quad\text{or}\quad b=0,
  \qquad
  \gcd\bigl(\gcd(d,e),|b-a|\bigr)=1.
\]
The depths \(a,b\) and outer degrees \(d,e\) are intrinsic valuation invariants of the
least output pair, not artifacts of a chosen factorization.

\subsection{Exact radical endpoint}

\tagLean The conjecture is equivalent to the localized-radical exclusion
\[
  u^\theta\notin\Z[1/3]
  \qquad\text{for every odd }u>1.
\]
It suffices to consider power-primitive odd \(u\).  For the canonical core \(w\), the
least index \(d\) above is also the exact degree of \(E=w^\theta\) over \(\Q\): the
binomial
\[
  X^d-\frac{A}{3^a}
\]
is irreducible and is the minimal polynomial of \(E\).  Thus a counterexample creates a
very special positive radical whose only rational-power denominator is supported at the
prime \(3\).

Taking logarithms gives the bilinear relation
\begin{equation}\label{eq:bilinear-logs}
  d\log w\log3=(\log A-a\log3)\log2.
\end{equation}
Baker's theory controls linear forms in logarithms with algebraic coefficients; it does not
exclude this product-of-logarithms relation.  The four-exponentials conjecture does.

\subsection{Three-smooth classification of auxiliary bases}

The formalized three-base determinant argument has the following exact consequence.

\begin{theorem}[Three-smooth-base theorem, repository theorem]\label{thm:three-smooth}
\tagLean Let \(x\in\Sol\setminus\Nzero\), and let \(B\in\N_{>0}\).  Then
\[
  B^x\in\Z
  \quad\Longleftrightarrow\quad
  B=2^u3^v\text{ for some }u,v\in\Nzero.
\]
\end{theorem}

The reverse implication is elementary.  The forward implication is a specialized
six-exponentials determinant proof: a prime divisor of \(B\) outside \(\{2,3\}\) supplies
a third multiplicatively independent base and forces \(x\) to be rational, hence integral.

\subsection{Counting and local progression obstructions}

\tagLean Exact coordinates give exact unit-block and dyadic-block counts.  In particular,
if failure occurs then
\begin{equation}\label{eq:total-asymptotic}
  N_\beta(T):=
  \#\{(n,k)\in\Nzero^2\setminus\{(0,0)\}:n+k\beta<T\}
  =\frac{T^2}{2\beta}+O_\beta(T).
\end{equation}
The repository proves the corresponding sharp limit, an equivalence between the conjecture
and subquadratic candidate growth below \(2^T\), and convergence of exact block averages
for every finite trigonometric polynomial.  The continuous-test and interval-indicator
extension is a paper-level Weyl consequence in the combined report.

\tagLean For every \(r\ge2\), set
\[
  c_r(\theta):=
  \left|\theta(\theta-1)\cdots(\theta-r+1)\right|.
\]
Candidates cannot occupy all points
\[
  n,n+h,\ldots,n+rh
\]
when
\begin{equation}\label{eq:higher-difference-condition}
  c_r(\theta)h^r<n^{r-\theta}.
\end{equation}
This is proved by an iterated mean-value theorem for forward differences: the \(r\)-th
difference is a nonzero integer but has absolute value below one.  We optimize the order
\(r\) in \cref{sec:finite-difference}.

\section{Attempts at a full proof and their exact failure points}\label{sec:attempts}

The following routes were pushed far enough to identify a precise obstruction.  Recording
the obstruction is important: it prevents a reformulation already compatible with the
exact monoid from being mistaken for a proof.

\subsection{Shrinking the three-base determinant to two bases}

The formalized determinant method has three independent logarithmic directions against two
exponent directions.  Its arithmetic lower bound and analytic interpolation upper bound
then have enough dimensional slack to force a contradiction.  With only \(\log2,\log3\),
the balance lands at the classical four-exponentials threshold.  Repackaging the same grid,
changing a few exponents, or using a larger square matrix does not create the missing
dimension.  A successful two-base determinant must use additional arithmetic information,
such as simultaneous non-Archimedean divisibility or the exact primitive normalization.

\subsection{The radical endpoint}

The reduction to an odd power-primitive \(w\) with \(w^\theta\in\overline\Q\), indeed with
an exact binomial minimal polynomial, is extremely sharp.  It does not itself contradict a
known theorem.  Even the first-looking special value
\[
  3^{\log_2 3}=e^{(\log3)^2/\log2}
\]
is not presently known to be irrational.  Gelfond--Schneider requires an algebraic
irrational exponent; here \(\log_2 3\) is transcendental.  Linear forms in logarithms do
not control the quadratic logarithmic expression.

\subsection{Equidistribution against exponentially shrinking targets}

A counterexample generates the irrational rotation orbit \(\{k\beta\}\), and the exact
block description makes that orbit equidistributed.  The integrality gap at height \(k\),
however, is exponentially small, of scale roughly \(M^{-k}\).  Ordinary discrepancy and
Dirichlet approximation control fixed or polynomially shrinking intervals, not this
exponential scale.  Equidistributed sequences with bounded partial quotients can avoid such
targets indefinitely.  A contradiction would require Diophantine information special to a
\(\beta\) satisfying \(2^\beta,3^\beta\in\N\), not generic rotation theory.

\subsection{Counting and additive combinatorics}

Under failure the candidate set has only \(O(\log X)\) points in \([X,2X)\).  This is far
below the density regimes of Roth, Szemer\'edi, or standard quantitative progression
results.  General \(S\)-unit theory can show finiteness of many projective additive
patterns, but the required conclusion is nonexistence.  The optimized finite-difference
analysis below makes the mismatch quantitative.

\subsection{Descent and perfect powers}

The least output pair has no common depth, common perfect-power degree, or matched affine
decomposition.  Yet the exact monoid predicts infinitely many later pairs with such
features.  Descent simply divides their coordinate pair by a common factor and returns to
another allowed point.  To close the problem, descent must produce a point below the
canonical least generator; the exact coordinate criteria in \cref{sec:perfect-powers}
show exactly when that can and cannot happen.

\subsection{A proof-audit criterion}

Any claimed proof must add an ingredient not already implied by
\(\Sol=\Nzero+\Nzero\beta\).  The following facts are all compatible with a
counterexample and cannot alone settle the conjecture:
\begin{itemize}
\item zero density or logarithmic-square counting;
\item equidistribution of fractional parts;
\item uniqueness of the primitive odd core;
\item transcendence of every noninteger solution;
\item arbitrary but finite computational exclusion;
\item fixed-order or variable-order local progression obstructions without a mechanism
forcing the required progression.
\end{itemize}

\section{Rational-function rigidity and nonlinear closure}\label{sec:rational-rigidity}

The exact generator theorem has a consequence not isolated in the starting report: a
hypothetical solution monoid is rigid against every nonlinear rational operation over the
algebraic numbers.

\subsection{The rational-function intersection theorem}

\begin{theorem}[Rational-function rigidity at the primitive generator]\label{thm:rational-rigidity}
\tagPaper Assume failure, and let \(\beta\) be the canonical generator from
\cref{thm:exact-monoid}.  Let \(R(T)\in\Qbar(T)\) have no pole at \(\beta\).  Then
\[
  \boxed{\quad
  R(\beta)\in\Sol
  \quad\Longleftrightarrow\quad
  R(T)=n+kT\text{ in }\Qbar(T)
  \text{ for some }n,k\in\Nzero.
  \quad}
\]
\end{theorem}

\begin{proof}
The reverse implication follows from additive closure:
\(R(\beta)=n+k\beta\in\Sol\).

For the forward implication, write \(R=P/Q\) with \(P,Q\in\Qbar[T]\) and
\(Q(\beta)\ne0\).  If \(R(\beta)\in\Sol\), exact coordinates give
\(R(\beta)=n+k\beta\) for some \(n,k\in\Nzero\).  Therefore
\[
  P(\beta)-(n+k\beta)Q(\beta)=0.
\]
By \cref{cor:beta-trans}, evaluation at \(\beta\) is injective on \(\Qbar[T]\).  Hence
\(P(T)-(n+kT)Q(T)\) is the zero polynomial, proving the claim.
\end{proof}

This theorem permits poles and cancellation and allows the full algebraic closure as
coefficient field.  It can be read as the exact intersection
\[
  \Sol\cap\Qbar(\beta)=\Nzero+\Nzero\beta,
\]
with the additional information that every rational expression realizing an element of
the intersection is identically the corresponding affine expression.

\begin{corollary}[Rational self-maps]\label{cor:rational-selfmaps}
\tagPaper Under failure, the rational functions over \(\Qbar\) that are regular on
\(\Sol\) and map \(\Sol\) into itself are exactly
\[
  T\longmapsto n+kT,
  \qquad n,k\in\Nzero.
\]
\end{corollary}

\begin{proof}
Such a function sends \(\beta\) to \(\Sol\), so
\cref{thm:rational-rigidity} identifies it.  Conversely, natural translations and natural
dilations preserve \(\Nzero+\Nzero\beta\).
\end{proof}

More generally, if \(x_i=n_i+k_i\beta\in\Sol\) and
\(F\in\Qbar(X_1,\ldots,X_r)\) is regular at the tuple, then
\[
  F(x_1,\ldots,x_r)\in\Sol
\]
exactly when the one-variable specialization
\[
  F(n_1+k_1T,\ldots,n_r+k_rT)
\]
is identically \(p+qT\) with \(p,q\in\Nzero\).  Thus every rational operation on finitely
many solutions is reduced to an elementary identity in one indeterminate.

\subsection{Exact multiplication criterion}

\begin{theorem}[Products of two solutions]\label{thm:product-criterion}
\tagPaper For all \(x,y\in\Sol\), unconditionally,
\[
  \boxed{\quad
  xy\in\Sol
  \quad\Longleftrightarrow\quad
  x\in\Nzero\ \text{or}\ y\in\Nzero.
  \quad}
\]
\end{theorem}

\begin{proof}
If one factor, say \(x\), is a nonnegative integer, then \(xy\) is a natural-number
multiple of a solution, so \(xy\in\Sol\).

For the converse, there is nothing to prove if the conjecture is true.  Under failure write
uniquely
\[
  x=n+k\beta,
  \qquad y=m+\ell\beta.
\]
If both are nonintegral, then \(k,\ell>0\), and
\[
  xy=nm+(n\ell+mk)\beta+k\ell\beta^2.
\]
Membership in \(\Sol\) would make this equal to \(r+s\beta\), producing a nonzero
quadratic polynomial over \(\Q\) that vanishes at the transcendental number \(\beta\).
Equivalently, apply \cref{thm:rational-rigidity} to
\((n+kT)(m+\ell T)\).  Both are impossible.
\end{proof}

\begin{corollary}[Finite products]\label{cor:finite-products}
\tagPaper Let \(x_1,\ldots,x_r\in\Sol\) be positive.  Then
\[
  \prod_{i=1}^r x_i\in\Sol
  \quad\Longleftrightarrow\quad
  \text{at most one }x_i\text{ is nonintegral}.
\]
If zero factors are allowed, a zero product is of course a solution.
\end{corollary}

\begin{proof}
Under failure, after removing positive integral factors, the product is a polynomial in
\(\beta\).  Its degree equals the number of noninteger factors and its leading coefficient
is positive.  A degree at least two is excluded by
\cref{thm:rational-rigidity}; degree zero or one is preserved by natural dilation.
\end{proof}

\begin{corollary}[Internal power test]\label{cor:power-test}
\tagPaper For \(x\in\Sol\) and every integer \(r\ge2\),
\[
  x^r\in\Sol\quad\Longleftrightarrow\quad x\in\Nzero.
\]
In particular, \(x\in\Nzero\) exactly when \(x^2\in\Sol\).
\end{corollary}

\subsection{Exact quotient criterion}

The same evaluation argument classifies division, which is more delicate because even the
true solution set \(\Nzero\) is not closed under arbitrary quotients.

\begin{theorem}[Quotients of solutions]\label{thm:quotient-criterion}
\tagPaper Assume failure and write
\[
  x=n+k\beta,\qquad y=m+\ell\beta>0.
\]
Then \(x/y\in\Sol\) if and only if exactly one of the following applicable conditions
holds:
\begin{enumerate}[label=\textup{(\alph*)}]
\item \(\ell=0\), so \(y=m\in\N_{>0}\), and \(m\mid n\) and \(m\mid k\);
\item \(\ell>0\), and \(x=qy\) for some \(q\in\Nzero\), equivalently
\(n=qm\) and \(k=q\ell\).
\end{enumerate}
In particular, if \(x,y\) are both noninteger, then
\[
  \frac{x}{y}\in\Sol
  \quad\Longleftrightarrow\quad
  \frac{x}{y}\in\Nzero.
\]
\end{theorem}

\begin{proof}
Suppose \(x/y=r+s\beta\) with \(r,s\in\Nzero\).  Clearing the denominator gives
\[
  n+k\beta
  =mr+(ms+\ell r)\beta+\ell s\beta^2.
\]
Transcendence forces \(\ell s=0\) and equality of the remaining coefficients.
If \(\ell=0\), then \(m>0\) and the equations are \(n=mr\), \(k=ms\).
If \(\ell>0\), then \(s=0\) and \(n=mr\), \(k=\ell r\), so \(x=ry\).
The converses follow by direct substitution.
\end{proof}

A useful special case is
\[
  x>0,
\frac1x\in\Sol
  \quad\Longleftrightarrow\quad x=1.
\]
Thus a hypothetical solution monoid is additively rich but almost maximally hostile to
multiplication, powers, reciprocals, and quotients.

\subsection{Closure formulations of the conjecture}

\begin{theorem}[Multiplicative closure equivalences]\label{thm:closure-equivalences}
\tagPaper The following statements are equivalent.
\begin{enumerate}[label=\textup{(\roman*)}]
\item The Alaoglu--Erd\H{o}s conjecture is true.
\item \(\Sol\) is closed under multiplication.
\item \(\Sol\) is closed under squaring.
\item For every real \(x\),
\[
  2^x,3^x\in\Z
  \quad\Longrightarrow\quad
  2^{x^2},3^{x^2}\in\Z.
\]
\end{enumerate}
\end{theorem}

\begin{proof}
If the conjecture holds, \(\Sol=\Nzero\), which has the stated closure properties.
Multiplicative closure implies square closure.  Conditions (iii) and (iv) are the same
statement, and square closure plus \cref{cor:power-test} forces every solution to be
integral.
\end{proof}

For positive reals define the commutative operation
\[
  m\star n
  :=2^{(\log_2m)(\log_2n)}
  =m^{\log_2n}=n^{\log_2m}.
\]
If \(m=2^x\) and \(n=2^y\), then \(m\star n=2^{xy}\).

\begin{corollary}[Candidate star-product]\label{cor:star-product}
\tagPaper For \(m,n\in\Cand\),
\[
  m\star n\in\Cand
  \quad\Longleftrightarrow\quad
  m\text{ is a power of }2\ \text{or}\ n\text{ is a power of }2.
\]
The conjecture is therefore equivalent both to closure of \(\Cand\) under \(\star\) and
to the apparently weaker diagonal condition
\[
  m\in\Cand\quad\Longrightarrow\quad m\star m\in\Cand.
\]
\end{corollary}

These are genuine equivalences, not a proof: there is presently no independent mechanism
forcing the nonlinear closure.  Their value is that they identify the missing property as
a single crisp failure of semiring-like structure.

\subsection{Lean formalization cost}

The multiplication, quotient, and power criteria should be inexpensive to formalize from
\lean{PrimitiveGenerator.lean} and \lean{Transcendence.lean}.  A direct route avoids a
general rational-function library:
\begin{enumerate}
\item split on \lean{AlaogluErdosConjecture};
\item in the failure branch obtain the canonical generator and unique coordinates;
\item expand the relevant equality as a polynomial of degree at most two in \(\beta\);
\item invoke transcendence to show all rational coefficients vanish;
\item discharge the remaining natural-number divisibility conditions with \lean{omega}.
\end{enumerate}
The full rational-function theorem needs evaluation injectivity over \(\Qbar\), but the
core multiplication theorem only needs the already formalized transcendence over \(\Q\).

\section{Three-smooth outputs and transcendental product exponents}\label{sec:smooth-output}

The multiplication theorem above uses the exact monoid and is therefore close to Lean.  A
second argument based on six exponentials gives more: when nonlinear closure fails, the
relevant exponential values are usually not merely nonintegral but transcendental.

\subsection{The two outputs cannot both be three-smooth}

\begin{theorem}[Three-smooth output test]\label{thm:output-smooth}
\tagPaper For every \(x\in\Sol\), the following are equivalent:
\begin{enumerate}[label=\textup{(\roman*)}]
\item \(x\in\Nzero\);
\item both \(2^x\) and \(3^x\) are three-smooth;
\item \(6^x\) is three-smooth.
\end{enumerate}
\end{theorem}

\begin{proof}
Condition (i) immediately implies (ii), and (ii) is equivalent to (iii) because
\(2^x,3^x\) are positive integers and a product is three-smooth exactly when both factors
are.

For (ii)\(\Rightarrow\)(i), write
\[
  2^x=2^a3^b,
  \qquad
  3^x=2^c3^d
\]
with \(a,b,c,d\in\Nzero\).  Taking logarithms gives
\[
  x=a+b\theta,
  \qquad
  x=d+\frac c\theta.
\]
Therefore
\[
  b\theta^2+(a-d)\theta-c=0.
\]
The repository proves \(\theta\) transcendental, so every rational coefficient vanishes.
Thus \(b=c=0\), \(a=d\), and \(x=a\in\Nzero\).
\end{proof}

This elementary-looking theorem uses the transcendence of \(\theta\), but not the exact
conditional generator.  It is unconditional and should be straightforward to add to the
Lean development.

\begin{corollary}[Primitive-core trichotomy]\label{cor:core-trichotomy}
\tagPaper In the simultaneous least-output normalization
\[
  M=2^aw^d,
  \qquad A=3^bv^e,
\]
at most one of \(w=3\) and \(v=2\) can hold.  Equivalently, exactly one of the following
three smoothness patterns occurs under failure:
\begin{center}
\begin{tabular}{@{}ccc@{}}
\toprule
\(M\) & \(A\) & primitive cores\\
\midrule
three-smooth & not three-smooth & \(w=3,\ v\ne2\)\\
not three-smooth & three-smooth & \(w\ne3,\ v=2\)\\
not three-smooth & not three-smooth & \(w\ne3,\ v\ne2\)\\
\bottomrule
\end{tabular}
\end{center}
The missing fourth pattern is impossible.
\end{corollary}

Indeed an odd, three-smooth, power-primitive integer greater than one is exactly \(3\), and
a three-free, three-smooth, power-primitive integer greater than one is exactly \(2\).

\subsection{A sharp third-base dichotomy from six exponentials}

We use a standard form of the six-exponentials theorem
\cite{Ramachandra1968II,Waldschmidt2000}: if \(u_1,u_2\) are linearly independent over
\(\Q\) and \(v_1,v_2,v_3\) are linearly independent over \(\Q\), then at least one of
the six numbers \(e^{u_iv_j}\) is transcendental.

\begin{lemma}[Third-base transcendence]\label{lem:third-base-trans}
\tagClassical Let \(y\in\Sol\setminus\Nzero\), and let \(B\in\N_{>0}\) be not
three-smooth.  Then \(B^y\) is transcendental.
\end{lemma}

\begin{proof}
The three logarithms \(\log2,\log3,\log B\) are linearly independent over \(\Q\): a
rational relation would give a multiplicative relation among \(2,3,B\), impossible because
\(B\) has a prime factor outside \(\{2,3\}\).  Also \(1,y\) are linearly independent
over \(\Q\), since a rational solution is integral.  Apply six exponentials.  Five of the
six values
\[
  2,\ 3,\ B,\ 2^y,\ 3^y,\ B^y
\]
are algebraic---indeed, the first five are integers.  Therefore \(B^y\) is
transcendental.
\end{proof}

Together with \cref{thm:three-smooth}, this gives an exact algebraicity trichotomy for
integer bases:
\[
  B^y\in\Qbar
  \quad\Longleftrightarrow\quad
  B^y\in\Z
  \quad\Longleftrightarrow\quad
  B\in\Smooth
  \qquad(y\in\Sol\setminus\Nzero).
\]
There is no algebraic-irrational middle case.

\subsection{Products of noninteger solutions}

\begin{theorem}[Product-exponent dichotomy]\label{thm:product-trans}
\tagClassical Let \(x,y\in\Sol\).
\begin{enumerate}[label=\textup{(\alph*)}]
\item If \(x\in\Nzero\) or \(y\in\Nzero\), then \(2^{xy},3^{xy}\in\N\).
\item If both \(x\) and \(y\) are nonintegral, then each of \(2^{xy}\) and
\(3^{xy}\) is either an integer or transcendental, at least one is transcendental, and
\[
  \boxed{\quad 6^{xy}\text{ is transcendental}.\quad}
\]
\end{enumerate}
\end{theorem}

\begin{proof}
Part (a) is natural-number dilation of a solution.  For part (b), put
\(M_x=2^x\) and \(A_x=3^x\).  Then
\[
  2^{xy}=M_x^y,
  \qquad
  3^{xy}=A_x^y.
\]
If \(M_x\) is three-smooth, \cref{thm:three-smooth} makes \(M_x^y\) an integer; if not,
\cref{lem:third-base-trans} makes it transcendental.  The same holds for \(A_x^y\).
By \cref{thm:output-smooth}, \(M_x,A_x\) cannot both be three-smooth, so at least one
value is transcendental.

Finally \(6^x=M_xA_x\) is an integer that is not three-smooth.  Applying
\cref{lem:third-base-trans} with base \(6^x\) and exponent \(y\) gives
\((6^x)^y=6^{xy}\) transcendental.
\end{proof}

\begin{corollary}[Diagonal transcendence]\label{cor:diagonal-trans}
\tagClassical If \(x\in\Sol\setminus\Nzero\), then \(6^{x^2}\) is transcendental,
and at least one of \(2^{x^2}\), \(3^{x^2}\) is transcendental.  More precisely,
\[
\begin{aligned}
  2^{x^2}\in\N &\Longleftrightarrow 2^x\in\Smooth,\\
  3^{x^2}\in\N &\Longleftrightarrow 3^x\in\Smooth,
\end{aligned}
\]
and failure of either condition makes the corresponding value transcendental.
\end{corollary}

\begin{corollary}[Mixed three-smooth bases]\label{cor:mixed-base-trans}
\tagClassical Let \(x,y\in\Sol\setminus\Nzero\), and let \(u,v\ge1\).  Then
\[
  (2^u3^v)^{xy}
\]
is transcendental.
\end{corollary}

\begin{proof}
The integer
\[
  (2^u3^v)^x=(2^x)^u(3^x)^v
\]
has every prime occurring in either output.  At least one output has a prime outside
\(\{2,3\}\), so this base is not three-smooth.  Apply
\cref{lem:third-base-trans} with exponent \(y\).
\end{proof}

\subsection{A global second-level trichotomy}

The prime-support theorem in \cref{sec:lattice} makes the preceding dichotomy uniform over
all noninteger solutions, not dependent on a particular pair.

\begin{theorem}[Uniform second-level pattern]\label{thm:uniform-second-level}
\tagClassical Assume failure and let \(M=2^\beta\), \(A=3^\beta\).  For every
\(x,y\in\Sol\setminus\Nzero\), exactly one of the following three global patterns holds:
\begin{center}
\begin{tabularx}{0.94\textwidth}{@{}cYY@{}}
\toprule
condition on \((M,A)\) & \(2^{xy}\) & \(3^{xy}\)\\
\midrule
\(M\in\Smooth,\ A\notin\Smooth\) & integer & transcendental\\
\(M\notin\Smooth,\ A\in\Smooth\) & transcendental & integer\\
\(M\notin\Smooth,\ A\notin\Smooth\) & transcendental & transcendental\\
\bottomrule
\end{tabularx}
\end{center}
The chosen pattern is independent of \(x\) and \(y\).
\end{theorem}

\begin{proof}
Write \(x=n+k\beta\) with \(k>0\).  Then
\[
  2^x=2^nM^k,
  \qquad 3^x=3^nA^k.
\]
Multiplication by a power of the distinguished smooth prime cannot remove an outside prime,
so \(2^x\) is three-smooth exactly when \(M\) is; similarly for \(A\).  Now apply
\cref{thm:product-trans}.  The impossible fourth pattern is excluded by
\cref{thm:output-smooth}.
\end{proof}

This theorem suggests a clean three-case research split.  The one-smooth-output cases reduce
to the exceptional primitive cores \(w=3\) or \(v=2\); the generic case has both second
iterates transcendental.

\subsection{Algebraicity formulations equivalent to the conjecture}

\begin{theorem}[Algebraicity tests]\label{thm:algebraicity-equivalences}
\tagClassical Each of the following is equivalent to the Alaoglu--Erd\H{o}s conjecture.
\begin{enumerate}[label=\textup{(\roman*)}]
\item For every \(x\in\Sol\), \(6^{x^2}\) is algebraic.
\item For every \(x\in\Sol\), both \(2^{x^2}\) and \(3^{x^2}\) are algebraic.
\item For every \(x,y\in\Sol\), \(6^{xy}\) is algebraic.
\item For one fixed pair \(u,v\ge1\), every \(x\in\Sol\) satisfies
\((2^u3^v)^{x^2}\in\Qbar\).
\item For every \(x\in\Sol\), the integer \(6^x\) is three-smooth.
\end{enumerate}
\end{theorem}

\begin{proof}
If the conjecture is true, all relevant exponents are nonnegative integers, and all listed
values are integers; condition (v) is also immediate.  Conversely, a noninteger solution
contradicts (i)--(iv) by \cref{cor:diagonal-trans,cor:mixed-base-trans}, and contradicts
(v) by \cref{thm:output-smooth}.
\end{proof}

These formulations replace integrality by the apparently weaker target of algebraicity.
Six exponentials makes them equivalent.  The exact missing step is now visible: one would
need an independent reason that some nonlinear iterate such as \(6^{x^2}\) is algebraic
from the original two integer values.

\subsection{Formalization boundary}

The output-smooth theorem and the integer/noninteger parts of the product trichotomy use
only the already formalized transcendence of \(\theta\), unique factorization, and
\cref{thm:three-smooth}.  The transcendence upgrade requires a general
six-exponentials theorem.  The current repository determinant specializes to integer
entries and proves the three-smooth classification.  A promising formal route is to
replace the determinant lower bound \(|D|\ge1\) by a number-field norm bound for a nonzero
algebraic determinant.  This would expose the same interpolation upper bound through an
``algebraic entries of bounded height'' interface and recover
\cref{lem:third-base-trans}.

\section{Common perfect-power content and exact affine descent}\label{sec:perfect-powers}

Assume failure throughout this section and fix the canonical generator \(\beta\), with
\[
  M=2^\beta,\qquad A=3^\beta.
\]
For \(x=n+k\beta\), exact coordinates give
\begin{equation}\label{eq:output-coordinates}
  2^x=2^nM^k,
  \qquad
  3^x=3^nA^k.
\end{equation}
The coordinate gcd is an intrinsic arithmetic invariant of this output pair.

\subsection{Simultaneous perfect powers}

\begin{theorem}[Exact common-power criterion]\label{thm:common-power}
\tagPaper Let \(x=n+k\beta\in\Sol\), and let \(r\ge1\).  The following are equivalent:
\begin{enumerate}[label=\textup{(\roman*)}]
\item both \(2^x\) and \(3^x\) are perfect \(r\)-th powers in \(\N\);
\item \(x/r\in\Sol\);
\item \(r\mid n\) and \(r\mid k\).
\end{enumerate}
\end{theorem}

\begin{proof}
If \(2^x=U^r\) and \(3^x=V^r\) with \(U,V>0\), uniqueness of positive real
\(r\)-th roots gives
\[
  2^{x/r}=U,
  \qquad 3^{x/r}=V,
\]
so \(x/r\in\Sol\).  The converse follows by raising the two integer values at \(x/r\)
to the \(r\)-th power.

Finally write \(x/r=p+q\beta\) in unique coordinates.  Then
\[
  n+k\beta=rp+rq\beta,
\]
so \(n=rp\), \(k=rq\).  This is exactly the pair of divisibility conditions, and the
reverse construction is immediate.
\end{proof}

For \(x>0\), define the \emph{common output power degree}
\[
  g(x):=\max\{r\ge1:2^x\text{ and }3^x\text{ are both perfect }r\text{-th powers}\}.
\]

\begin{corollary}[Coordinate gcd]\label{cor:gcd-degree}
\tagPaper If \(x=n+k\beta>0\), then
\[
  \boxed{\quad g(x)=\gcd(n,k).\quad}
\]
The output pair is simultaneously power-primitive exactly when \(\gcd(n,k)=1\).
\end{corollary}

Thus every nonzero output pair has a unique decomposition
\[
  (2^x,3^x)=(U^g,V^g),
  \qquad g=\gcd(n,k),
\]
where \((U,V)\) is a simultaneously power-primitive output pair belonging to the solution
\(x/g\).  The repository proves affine primitivity for the least pair; this corollary
extends the intrinsic power degree to every point of the solution monoid.

\subsection{Matched affine decompositions}

A useful descent pattern strips the same distinguished-base depth and the same outer power
from both outputs.

\begin{theorem}[Matched affine decomposition criterion]\label{thm:affine-decomp}
\tagPaper Let \(x=n+k\beta\in\Sol\), \(d\ge1\), and \(r\in\Nzero\).  The following
are equivalent:
\begin{enumerate}[label=\textup{(\roman*)}]
\item there exist \(U,V\in\N_{>0}\) such that
\[
  2^x=2^rU^d,
  \qquad 3^x=3^rV^d;
\]
\item \((x-r)/d\in\Sol\);
\item \(0\le r\le n\), \(d\mid k\), and \(d\mid(n-r)\).
\end{enumerate}
When these conditions hold,
\[
  \frac{x-r}{d}=\frac{n-r}{d}+\frac{k}{d}\beta,
  \qquad
  U=2^{(x-r)/d},\quad V=3^{(x-r)/d}.
\]
\end{theorem}

\begin{proof}
Conditions (i) and (ii) are equivalent by positive-root uniqueness after removing the common
factor \(2^r\) or \(3^r\).  Expanding (ii) in unique coordinates gives
\[
  n+k\beta=r+d(p+q\beta)
\]
for some \(p,q\in\Nzero\).  Thus \(n=r+dp\) and \(k=dq\), which is exactly (iii).
The same equations give the formulas and the converse.
\end{proof}

For a noninteger solution \(x=n+k\beta\) with \(k>0\), Euclidean division
\(n=qk+r\), \(0\le r<k\), gives a canonical reduced affine presentation
\begin{equation}\label{eq:euclidean-affine}
  \boxed{\quad x=r+k(q+\beta),\quad}
\end{equation}
and hence
\[
  2^x=2^r(2^qM)^k,
  \qquad
  3^x=3^r(3^qA)^k.
\]
The normalization \(0\le r<k\) makes the presentation unique.  It separates the outer
affine degree \(k\) from the reduced common base depth \(r\).

\subsection{Why this does not yet descend below \texorpdfstring{\(\beta\)}{beta}}

If \(n,k\) have a common divisor, \cref{thm:common-power} descends to
\(x/\gcd(n,k)\).  If a matched decomposition exists, \cref{thm:affine-decomp} descends to
\((x-r)/d\).  But exact coordinates ensure that the descended point remains
\(p+q\beta\) with \(p,q\ge0\); it is below \(\beta\) only when \(q=0\), in which case
it is integral.  The least noninteger point \(\beta\) is already primitive with respect
to every such descent.  A successful descent argument therefore needs a transformation
not captured by common powers and common distinguished-base depths.

\section{Primitive-output statistics}\label{sec:primitive-counting}

The common-power criterion translates ordinary coprimality statistics of lattice points
into exact statistics of hypothetical integer-output pairs.

\subsection{Total and primitive counts}

For \(T>0\), let
\[
  N_\beta(T)
  =\#\{(n,k)\in\Nzero^2\setminus\{(0,0)\}:n+k\beta<T\}
\]
and
\[
  P_\beta(T)
  =\#\{(n,k)\in\Nzero^2:n+k\beta<T,\ \gcd(n,k)=1\}.
\]
The first counts nonzero solutions below \(T\); the second counts output pairs with no
nontrivial common perfect-power degree.  Recall
\[
  N_\beta(T)=\frac{T^2}{2\beta}+O_\beta(T).
\]

\begin{theorem}[Primitive output pairs]\label{thm:primitive-asymptotic}
\tagPaper As \(T\to\infty\),
\[
  \boxed{\quad
  P_\beta(T)
  =\frac{T^2}{2\beta\zeta(2)}+O_\beta(T\log T)
  =\frac{3T^2}{\pi^2\beta}+O_\beta(T\log T).
  \quad}
\]
Consequently,
\[
  \frac{P_\beta(T)}{N_\beta(T)}\longrightarrow\frac1{\zeta(2)}=\frac6{\pi^2}.
\]
\end{theorem}

\begin{proof}
Every nonzero coordinate pair has the unique factorization
\[
  (n,k)=d(n_0,k_0),
  \qquad d=\gcd(n,k),\quad \gcd(n_0,k_0)=1.
\]
Hence
\[
  N_\beta(T)=\sum_{d\ge1}P_\beta(T/d),
\]
and M\"obius inversion gives
\[
  P_\beta(T)=\sum_{d\le T}\mu(d)N_\beta(T/d).
\]
The lattice estimate is uniform in the form
\(N_\beta(U)=U^2/(2\beta)+O_\beta(U+1)\).  Therefore
\[
\begin{aligned}
  P_\beta(T)
  &=\frac{T^2}{2\beta}\sum_{d\le T}\frac{\mu(d)}{d^2}
    +O_\beta\!\left(\sum_{d\le T}\left(\frac Td+1\right)\right)\\
  &=\frac{T^2}{2\beta\zeta(2)}+O_\beta(T\log T).
\end{aligned}
\]
Divide by the total asymptotic to obtain the limit.
\end{proof}

The universal constant \(6/\pi^2\) is independent of the hypothetical generator.
Geometry contributes \(1/(2\beta)\); primitive lattice points contribute
\(1/\zeta(2)\).

\subsection{Distribution of maximal common power degree}

For fixed \(g\ge1\), let
\[
  P_{\beta,g}(T)
  :=\#\{x=n+k\beta:0<x<T,\ \gcd(n,k)=g\}.
\]

\begin{corollary}[Exact-degree distribution]\label{cor:degree-distribution}
\tagPaper For fixed \(g\ge1\),
\[
  P_{\beta,g}(T)=P_\beta(T/g)
  =\frac{3T^2}{\pi^2\beta g^2}
   +O_\beta\!\left(\frac Tg\log\frac{2T}{g}\right).
\]
The limiting probability that an output pair has maximal common power degree exactly
\(g\) is
\[
  \boxed{\quad \frac{6}{\pi^2g^2}.\quad}
\]
\end{corollary}

The probabilities sum to one.  Likewise, the two outputs are both perfect \(r\)-th powers
exactly when \(r\mid n,k\), so
\begin{equation}\label{eq:rth-power-frequency}
  \#\{0<x<T:2^x,3^x\text{ both perfect }r\text{-th powers}\}
  =N_\beta(T/r),
\end{equation}
and the limiting proportion is \(1/r^2\).

\subsection{Exact generating series}

The coordinate monoid has the Laplace series
\begin{equation}\label{eq:laplace-series}
\begin{aligned}
  Z_\beta(s)
  &:=\sum_{x\in\Sol\setminus\{0\}}e^{-sx}\\
  &=\sum_{n,k\ge0}e^{-s(n+k\beta)}-1\\
  &=\boxed{\frac1{(1-e^{-s})(1-e^{-\beta s})}-1},
  \qquad s>0.
\end{aligned}
\end{equation}
The primitive-output series is
\[
  Z_\beta^{\rm prim}(s)=\sum_{d\ge1}\mu(d)Z_\beta(ds).
\]
At the candidate level one similarly has the Dirichlet series
\[
  \sum_{m\in\Cand}m^{-s}
  =\frac1{(1-2^{-s})(1-M^{-s})},
  \qquad \Re s>0.
\]
These identities make the two free generators visible and give clean starting points for
weighted Tauberian refinements.

\subsection{Interpretation}

Roughly \(60.79\%\) of hypothetical output pairs would be simultaneously
power-primitive.  This is not pressure toward a contradiction: it is exactly the expected
primitive-lattice proportion for a free rank-two monoid.  Its value is diagnostic.  Any
descent argument must act on the primitive majority, not merely on the exceptional pairs
with a common outer power.

\section{Candidate divisibility and ambient root saturation}\label{sec:lattice}

The candidate monoid also has a simple valuation geometry.  It yields closure properties
not visible from the defining condition \(m^\theta\in\N\).

\subsection{A two-dimensional valuation cone}

Under failure write
\[
  M=2^aU,
  \qquad U>1\text{ odd}.
\]
A candidate with coordinates \((n,k)\) is
\[
  m(n,k)=2^nM^k=2^{n+ak}U^k.
\]
Set \(s=v_2(m)=n+ak\).  The map
\[
  m(n,k)\longmapsto(s,k)
\]
is a bijection from \(\Cand\) to
\[
  \Lambda_a:=\{(s,k)\in\Nzero^2:s\ge ak\}.
\]
Because \(U\) is odd and greater than one,
\begin{equation}\label{eq:divisibility-coordinates}
  m(n,k)\mid m(n',k')
  \quad\Longleftrightarrow\quad
  s\le s'\text{ and }k\le k'.
\end{equation}
Thus candidate divisibility is coordinatewise order on a lattice cone.

\begin{theorem}[Candidate gcd/lcm lattice]\label{thm:gcd-lcm}
\tagPaper Unconditionally, \(\Cand\) is closed under gcd and lcm.  Under failure, if
\(m_i=m(n_i,k_i)\) and \(s_i=n_i+ak_i\), then
\[
\begin{aligned}
  \gcd(m_1,m_2)
  &=m\!\left(\min(s_1,s_2)-a\min(k_1,k_2),\ \min(k_1,k_2)\right),\\
  \operatorname{lcm}(m_1,m_2)
  &=m\!\left(\max(s_1,s_2)-a\max(k_1,k_2),\ \max(k_1,k_2)\right).
\end{aligned}
\]
Candidate divisibility is therefore a distributive lattice.
\end{theorem}

\begin{proof}
If the conjecture is true, candidates are powers of two.  Under failure, gcd and lcm take
coordinatewise minima and maxima of the prime-valuation vectors of \(2^sU^k\).  The
resulting points remain in \(\Lambda_a\); coordinatewise min and max form a distributive
lattice.
\end{proof}

\begin{corollary}[Candidate divisors]\label{cor:candidate-divisors}
\tagPaper The number of candidate divisors of \(m(n,k)\), including \(1\) and the number
itself, is
\[
  \boxed{\quad
  (k+1)(n+1)+\frac{ak(k+1)}2.
  \quad}
\]
\end{corollary}

\begin{proof}
A candidate divisor corresponds to \((s',j)\in\Lambda_a\) with
\(0\le j\le k\) and \(aj\le s'\le n+ak\).  There are \(n+ak-aj+1\) choices for
\(s'\).  Sum this arithmetic progression over \(j=0,\ldots,k\).
\end{proof}

The gcd/lcm closure is compatible with failure, unlike the nonlinear star-product closure.
It is therefore structural information rather than a direct proof route.

\subsection{Fixed prime support}

Use the simultaneous normalization
\[
  M=2^aw^d,
  \qquad A=3^bv^e.
\]
For \(x=n+k\beta\),
\[
  2^x=2^{n+ak}w^{dk},
  \qquad
  3^x=3^{n+bk}v^{ek}.
\]

\begin{corollary}[Prime-support rigidity]\label{cor:prime-support}
\tagPaper Every noninteger solution \((k>0)\) has exactly the same odd prime support in
its base-\(2\) output, namely \(\supp(w)\), and exactly the same prime support away from
\(3\) in its base-\(3\) output, namely \(\supp(v)\).  The only support variation is
whether the distinguished primes \(2\) and \(3\) occur.
\end{corollary}

A counterexample would therefore lock the entire noninteger branch onto one fixed pair of
finite prime sets.  Computational searches should exploit this by enumerating primitive
supports or cores, not raw integers independently.

\subsection{Ambient root saturation and the index \texorpdfstring{\(d\)}{d}}

Use the canonical form \(M=2^aw^d\), where \(w>1\) is odd and the gcd of its positive
prime valuations is one.  Reparameterizing by the exponent of \(w\) gives
\begin{equation}\label{eq:candidate-cone}
  \Cand
  =\{2^sw^t:s,t\in\Nzero,\ d\mid t,\ ds\ge at\}.
\end{equation}

For a multiplicative monoid \(H\subseteq\N_{>0}\), define its \emph{ambient root
saturation}
\[
  \Sat(H):=\{u\in\N_{>0}:u^r\in H\text{ for some }r\ge1\}.
\]

\begin{theorem}[Exact ambient root saturation]\label{thm:root-saturation}
\tagPaper Under failure,
\[
  \boxed{\quad
  \Sat(\Cand)=\{2^sw^t:s,t\in\Nzero,\ ds\ge at\}.
  \quad}
\]
For \(u=2^sw^t\) in the saturation, the least positive \(r\) for which
\(u^r\in\Cand\) is
\[
  \boxed{\quad r_{\min}(u)=\frac d{\gcd(d,t)}.\quad}
\]
\end{theorem}

\begin{proof}
Suppose first that \(u=2^sw^t\) and \(ds\ge at\).  Put
\(r=d/\gcd(d,t)\).  Then \(d\mid rt\), and scaling the cone inequality gives
\(d(rs)\ge a(rt)\).  Hence \(u^r\in\Cand\).  The divisibility condition
\(d\mid rt\) proves minimality.

Conversely suppose \(u^r\in\Cand\), so
\[
  u^r=2^Sw^{dk}
\]
with \(dS\ge adk\).  No prime outside \(\{2\}\cup\supp(w)\) divides \(u\).  Write
\(c_p=v_p(w)\).  For every \(p\mid w\),
\[
  r v_p(u)=dkc_p.
\]
Since \(\gcd\{c_p:p\mid w\}=1\), B\'ezout gives \(dk/r\in\Z\); call it \(t\).
The odd part of \(u\) is then \(w^t\), while \(S=rs\) with \(s=v_2(u)\).  Substitution
in the cone inequality yields \(ds\ge at\).
\end{proof}

\begin{corollary}[Ambient saturation index]\label{cor:ambient-index}
\tagPaper The candidate monoid equals its ambient root saturation exactly when \(d=1\).
More generally,
\[
  d=\min\{r\ge1:\exists s\in\Nzero,\ (2^sw)^r\in\Cand\}.
\]
\end{corollary}

A terminology caution is important.  In standard affine-semigroup theory, normality is
measured relative to the group generated by the semigroup.  Here that group already has
\(t\equiv0\pmod d\), and \(\Cand\) is normal in that intrinsic lattice for every \(d\).
The theorem instead measures saturation in the larger ambient lattice containing every
\(2^sw^t\).  Thus \(d\) is a finite-index \emph{over-lattice saturation defect}, not a
failure of intrinsic semigroup normality.  This is precisely parallel to its field-theoretic
meaning as the least rational-power index and exact radical degree.

\section{Optimizing the higher-difference obstruction}\label{sec:finite-difference}

The arbitrary-order theorem in \lean{HigherDifference.lean} permits the order to grow with
the initial candidate size.  Optimizing that order strengthens every fixed-order corollary,
but also quantifies why the method does not presently contradict the conditional monoid.

\subsection{The exact derivative coefficient}

For \(r\ge2\) and \(1<\theta<2\), set
\[
  c_r=\left|\theta(\theta-1)\cdots(\theta-r+1)\right|.
\]
Because the factors from \(2-\theta\) onward are positive after taking absolute values,
\begin{equation}\label{eq:cr-gamma}
  c_r
  =\frac{\theta(\theta-1)}{\Gamma(2-\theta)}\Gamma(r-\theta).
\end{equation}
Stirling's formula gives
\begin{equation}\label{eq:cr-root}
  c_r^{1/r}
  =\frac re\left(1+O\!\left(\frac{\log r}{r}\right)\right).
\end{equation}

The formalized progression theorem applies whenever
\[
  h<c_r^{-1/r}n^{1-\theta/r}.
\]
Using \eqref{eq:cr-root}, the right side is
\begin{equation}\label{eq:threshold-approx}
  \frac{en}{r}
  \exp\!\left(-\frac{\theta\log n}{r}\right)
  \left(1+O\!\left(\frac{\log r}{r}\right)\right).
\end{equation}
Ignoring the lower-order factor, the logarithmic derivative with respect to \(r\) is
\[
  -\frac1r+\frac{\theta\log n}{r^2},
\]
so the maximum occurs at \(r\sim\theta\log n\).  At that value the explicit factor \(e\)
and the exponential \(e^{-1}\) cancel.

\begin{theorem}[Optimized variable-order obstruction]\label{thm:optimized-difference}
\tagPaper Let \(0<\varepsilon<1\).  For all sufficiently large \(n\), choose an integer
\[
  r_n=\theta\log n+O(1),\qquad r_n\ge2.
\]
Then candidates cannot occupy all terms of
\[
  n,n+h,\ldots,n+r_nh
\]
whenever
\[
  1\le h\le(1-\varepsilon)\frac{n}{\theta\log n}.
\]
Moreover, the largest leading-order step scale obtainable by optimizing the existing
curvature condition over \(r\) is
\[
  \boxed{\quad h_{\max}(n)\sim\frac{n}{\theta\log n}.\quad}
\]
\end{theorem}

\begin{proof}
Choose \(r_n\) to be either nearest integer to \(\theta\log n\).  Then \(r_n\to\infty\),
so \eqref{eq:threshold-approx} applies and its quotient by
\(n/(\theta\log n)\) tends to one.  The fixed \((1-\varepsilon)\) margin absorbs the
Stirling error and the \(O(1)\) integer rounding.  The formalized condition
\eqref{eq:higher-difference-condition} then excludes the progression.

For the optimality assertion, the leading approximation in
\eqref{eq:threshold-approx} is a unimodal function of positive real \(r\), with its unique
maximum at \(\theta\log n\).  Regimes where \(r\) remains bounded or grows outside a
constant multiple of \(\log n\) give a strictly smaller order or leading constant.
\end{proof}

\subsection{Comparison with conditional candidate spacing}

A dyadic block \([X,2X)\) is an interval of length one in exponent space.  Exact coordinates
give
\begin{equation}\label{eq:dyadic-count}
  \#(\Cand\cap[X,2X))
  =\frac{\log X}{\beta\log2}+O_\beta(1).
\end{equation}
The average arithmetic spacing in the block is therefore
\[
  \sim\beta\log2\,\frac{X}{\log X}.
\]
At \(n\asymp X\), the optimized curvature threshold is
\[
  \sim\frac1\theta\frac{X}{\log X}.
\]
The ratio is
\begin{equation}\label{eq:spacing-ratio}
  \frac{\beta\log2}{1/\theta}
  =\boxed{\beta\log3}.
\end{equation}
The kernel-verified bound \(\beta\ge12\) already makes this ratio greater than
\(13.18\); the software-level \(\beta>26\) bound makes it greater than \(28.56\).

The same constant appears in a cardinality comparison.  The forbidden progression would
require
\[
  r_n+1\sim\theta\log X
\]
candidates in one structured configuration, while the entire dyadic block contains only
\[
  \sim\frac{\log X}{\beta\log2}
\]
candidates.  Their ratio is again \(\beta\log3\).

\begin{remark}[Quantified barrier]\label{rem:difference-barrier}
The variable-order optimization materially strengthens the local theorem, but total
candidate density is too small by a fixed factor of at least \(\beta\log3\).  No argument
using only the number of candidates in a dyadic block can force the progression required
by the current finite-difference functional.  A successful continuation must exploit
additional arithmetic correlation---for example, valuation congruences that force a
special progression---or find a different integer-valued functional with a better
curvature-to-spacing constant.
\end{remark}

This is a limitation of the method, not evidence against the conjecture.  Its value is to
replace the vague statement ``density seems too weak'' by a sharp asymptotic constant and
a concrete target for any improved functional.

\section{A directed-rounding exclusion through \texorpdfstring{\(2^{26}\)}{2 to the 26}}\label{sec:computation}

This section records a new finite computation and its exact trust boundary.  It is not a
substitute for Lean: its purpose is to strengthen the explicit zero-free region, test the
conditional models, and provide data for future certificate generation.

\subsection{What must be certified for each integer}

For a non-power of two \(m\), let
\[
  y=m^\theta,
  \qquad \theta=\frac{\log3}{\log2}.
\]
To prove that \(y\notin\N\), it suffices to find \(n\in\N\) and certify the strict pair
\[
  n<m^\theta<n+1.
\]
Avoiding direct exponentiation, these are equivalent to
\begin{align}
  \log m\,\log3-\log n\,\log2 &>0,
    \label{eq:lower-log-gap}\\
  \log(n+1)\,\log2-\log m\,\log3 &>0.
    \label{eq:upper-log-gap}
\end{align}
All logarithms and products are positive in the scanned range, so interval directions are
simple and monotone.

The verifier uses ordinary \texttt{long double} only to propose
\(n\approx\lfloor m^\theta\rfloor\).  It tests the offsets
\[
  0,-1,1,-2,2,-3,3.
\]
A proposal is accepted only when both \eqref{eq:lower-log-gap} and
\eqref{eq:upper-log-gap} have a strictly positive \emph{directed lower bound} computed by
MPFR.  Failure to locate an interval is reported as a failure, never silently classified as
nonintegral.  Thus locator error can create a false alarm but cannot create a false proof.

\subsection{Directed enclosures}

At precision \(p=192\) bits, the program constructs
\[
\begin{aligned}
  L_m^-&\le\log m\le L_m^+, &
  L_2^-&\le\log2\le L_2^+, &
  L_3^-&\le\log3\le L_3^+,\\
  L_n^+&\ge\log n, &&
  L_{n+1}^-&\le\log(n+1).&&
\end{aligned}
\]
It then rounds downward in
\[
  \delta_-=L_m^-L_3^- - L_n^+L_2^+,
  \qquad
  \delta_+=L_{n+1}^-L_2^- - L_m^+L_3^+.
\]
Thus \(\delta_->0\) proves \eqref{eq:lower-log-gap}, and \(\delta_+>0\) proves
\eqref{eq:upper-log-gap}.  The input integers are below \(2^{42}\) on the \(n\)-side and
are exactly representable at 192-bit precision.

Powers of two are skipped, not tested: if \(m=2^q\), then
\(m^\theta=3^q\) is the known integral candidate.  Every other positive integer is tested.
An exact integral value at a non-power of two could satisfy neither strict interval for any
\(n\), so it would necessarily be reported as a failure.

\subsection{Execution environment and source identity}

The recorded run used:
\begin{itemize}
\item MPFR \(4.2.2\), linked as \texttt{libmpfr.so.6};
\item GCC \(14.2.0\) with \texttt{-O3 -fopenmp};
\item GNU OpenMP with five worker threads;
\item Linux x86-64;
\item 192-bit MPFR precision and directed modes \texttt{MPFR\_RNDD}/\texttt{MPFR\_RNDU}.
\end{itemize}
The exact executed C source is reproduced in \cref{app:scanner}.  Its SHA-256 digest is
\begin{center}
\small\ttfamily
f888c0d89be9366f01e2d26f339dcbe446556d39d4c5dbcf8f579ec0d0052e1d
\end{center}
The execution image contained the MPFR shared library but not the development header, so
the source has a guarded minimal declaration of the MPFR~4.x ABI.  On an ordinary system
with \texttt{mpfr.h}, the standard header branch is used.  This ABI detail is part of the
software trust boundary and is disclosed rather than hidden.

\subsection{Result}

There are \(2^{26}-1=67{,}108{,}863\) positive integers below \(2^{26}\).  Exactly
\(26\) are powers of two, from \(2^0\) through \(2^{25}\).  Hence the verifier tests
\[
  67{,}108{,}863-26
  =\boxed{67{,}108{,}837}
\]
non-powers.  The range was divided at multiples of \(2^{22}\) into one initial chunk
\([1,2^{22})\) and fifteen chunks of width \(2^{22}\); all chunks
returned zero failures.  The complete chunk audit appears in \cref{app:chunk-audit}.

\begin{theorem}[Software-level finite exclusion]\label{thm:finite-26}
\tagComp For every integer \(m\) with \(1\le m<2^{26}\),
\[
  m^\theta\in\N
  \quad\Longrightarrow\quad
  m\text{ is a power of }2.
\]
Consequently, at the stated software trust level,
\[
  x\in\Sol\setminus\Nzero
  \quad\Longrightarrow\quad
  \boxed{x>26}.
\]
\end{theorem}

\begin{proof}[Explanation of the consequence]
The verifier establishes the first implication for every non-power input; powers of two are
the known candidates.  If \(x\) were nonintegral, then \(m=2^x\) would be an integer
candidate but not a power of two.  The finite exclusion gives \(m\ge2^{26}\).  Equality is
itself a power of two and would force \(x=26\), so in fact \(m>2^{26}\), hence \(x>26\).
\end{proof}

\subsection{Closest approaches}

The smallest certified lower logarithmic margin over the full scan is
\[
  3.6447144482202515\times10^{-20}
\]
at
\[
  (m,n)=(58{,}570{,}438,\ 2{,}048{,}703{,}434{,}093).
\]
A separate 100-decimal evaluation gives
\[
  m^\theta-n
  =1.07725158750919017548\ldots\times10^{-7}.
\]

The smallest certified upper logarithmic margin is
\[
  5.4132120119800561\times10^{-21}
\]
at
\[
  (m,n)=(29{,}411{,}704,\ 687{,}581{,}893{,}443),
\]
where
\[
  n+1-m^\theta
  =5.36974926710988874953\ldots\times10^{-9}.
\]
The same logarithmic upper margin recurs at the scaled point
\((58{,}823{,}408,2{,}062{,}745{,}680{,}331)\), as expected from multiplication of
\(m\) by two and the identity \((2m)^\theta=3m^\theta\).

Both extremal inequalities were independently reevaluated at 100 decimal digits.  This
secondary check is not an independent proof system, but it is a useful guard against
transcription and aggregation errors.  The certified margins are many orders of magnitude
larger than the nominal 192-bit rounding scale of the logarithmic operations.

\subsection{Trust boundary and route to a kernel theorem}

The computation trusts the C source, GCC code generation, OpenMP scheduling, the MPFR and
GMP shared libraries, the x86-64 arithmetic environment, and the operating system.  MPFR's
correct-rounding contract makes the numerical reasoning crisp, but it does not remove
those software assumptions.

The natural next step is certificate extraction rather than a still larger blind scan.
For each block, one can record rational lower and upper approximants to \(\theta\) and
integer inequalities sufficient to replay every exclusion.  Compression should exploit
monotonicity and convexity so that one certificate covers a run of consecutive \(m\),
rather than storing one record per input.  The existing \(4096\) Lean proof provides the
checker pattern.  A practical intermediate target is \(2^{16}\), followed by a generated
hierarchical certificate through \(2^{26}\).

\section{A consolidated catalogue of equivalent formulations}\label{sec:equivalences}

The repository already contains many equivalences: candidate powers of two, odd localized
radicals, primitive radicals, least-generator nonexistence, and subquadratic counting.  The
new results add nonlinear closure and algebraicity tests.  The following theorem collects
the most useful forms in one place.

\begin{theorem}[Equivalent targets]\label{thm:equivalent-targets}
The following assertions are equivalent.
\begin{enumerate}[label=\textup{(\arabic*)}]
\item \tagLean The Alaoglu--Erd\H{o}s conjecture: \(\Sol=\Nzero\).
\item \tagLean Every candidate is a power of two.
\item \tagLean There is no odd \(u>1\) with \(u^\theta\in\Z[1/3]\); equivalently, none
with \(u\) power-primitive.
\item \tagPaper \(\Sol\) is closed under multiplication.
\item \tagPaper \(\Sol\) is closed under squaring.
\item \tagPaper \(\Cand\) is closed under the operation
\(m\star n=2^{\log_2m\log_2n}\).
\item \tagPaper Every \(m\in\Cand\) satisfies \(m\star m\in\Cand\).
\item \tagPaper Every \(x\in\Sol\) has both outputs \(2^x,3^x\) three-smooth.
\item \tagPaper Every \(x\in\Sol\) has \(6^x\) three-smooth.
\item \tagClassical Every \(x\in\Sol\) has \(6^{x^2}\in\Qbar\).
\item \tagClassical Every \(x\in\Sol\) has both \(2^{x^2},3^{x^2}\in\Qbar\).
\item \tagClassical Every \(x,y\in\Sol\) has \(6^{xy}\in\Qbar\).
\item \tagLean The candidate count below \(2^T\) is \(o(T^2)\).
\end{enumerate}
\end{theorem}

\begin{proof}
The repository proves the equivalence of (1)--(3) and (13).  The equivalence of (1) with
(4)--(7) is \cref{thm:closure-equivalences,cor:star-product}.  The equivalence with
(8)--(9) is \cref{thm:output-smooth}.  The equivalence with (10)--(12) is
\cref{thm:algebraicity-equivalences}.
\end{proof}

The formulations are not equally promising.  Counting is very far from the final
integrality statement; multiplication closure is exact but has no known independent
source; the radical formulation is closest to established transcendence theory; and the
three-smooth/algebraicity formulations offer a possible route through a third-base theorem.

\subsection{A third-base manufacturing principle}

\Cref{thm:three-smooth} suggests a particularly concrete target.  Given a hypothetical
noninteger \(\beta\), it would suffice to construct \emph{one} positive integer
\(B\notin\Smooth\) such that \(B^\beta\in\Z\).  This would immediately contradict the
formalized three-smooth-base theorem.

The obvious choices are \(B=M=2^\beta\), \(B=A=3^\beta\), or \(B=MA=6^\beta\).
Their \(\beta\)-th powers are
\[
  2^{\beta^2},\qquad 3^{\beta^2},\qquad 6^{\beta^2}.
\]
Thus ``manufacture a third base'' is exactly the square-closure problem in a different
language.  The six-exponentials result makes the stakes sharper: whenever the chosen base
is not three-smooth, its \(\beta\)-th power is already known to be transcendental.  Any
proof of algebraicity would therefore finish the conjecture at once.

This framing is useful because it isolates the missing bridge:
\[
  \boxed{\text{derive one nonlinear algebraic value from the two linear integer values}.}
\]
Neither unique factorization, ordinary exponent laws, nor linear forms in logarithms
currently supply that bridge.

\section{Prioritized future research program}\label{sec:research}

The following agenda is ordered by a combination of mathematical leverage, falsifiability,
and formalization value.  It is intentionally technical rather than promotional.

\subsection{Priority 1: formalize the new elementary rigidity theorems}

\tagProposal Add Lean modules for:
\begin{itemize}
\item three-smooth output equivalence, \cref{thm:output-smooth};
\item exact multiplication and power criteria,
\cref{thm:product-criterion,cor:power-test};
\item quotient classification, \cref{thm:quotient-criterion};
\item common perfect-power and affine-decomposition criteria,
\cref{thm:common-power,thm:affine-decomp};
\item candidate gcd/lcm closure and divisor enumeration,
\cref{thm:gcd-lcm,cor:candidate-divisors};
\item ambient root saturation, with the intrinsic-lattice terminology caveat;
\item primitive lattice asymptotics, initially as exact finite M\"obius identities and later
through Mathlib asymptotics.
\end{itemize}

These proofs are short relative to the existing determinant and Gelfond--Schneider
machinery.  Formalizing them would give the project a larger library of reusable
``counterexample-world eliminators'' and guard against subtle sign, divisibility, or
normalization errors.

A suggested source split is:
\begin{center}
\begin{tabularx}{0.95\textwidth}{@{}lY@{}}
\toprule
module & contents\\
\midrule
\texttt{RationalOperationRigidity.lean} & products, powers, quotients, specialized polynomial evaluation\\
\texttt{OutputSmoothness.lean} & three-smooth output test and primitive-core trichotomy\\
\texttt{CommonPowerContent.lean} & exact common power degree and affine descent\\
\texttt{CandidateDivisibility.lean} & gcd/lcm lattice, candidate divisors, support rigidity\\
\texttt{AmbientRootSaturation.lean} & exact over-lattice saturation and least root exponent\\
\texttt{PrimitiveOutputCounting.lean} & finite M\"obius identities and asymptotic wrappers\\
\bottomrule
\end{tabularx}
\end{center}

\subsection{Priority 2: turn the \texorpdfstring{\(2^{26}\)}{2-to-the-26} scan into replayable certificates}

\tagProposal The next computational objective should not simply be \(2^{27}\).  It should
be a compact certificate format accepted by a small Lean checker.  Useful compression
mechanisms include:
\begin{enumerate}
\item rational bounds \(p/q<\theta<r/s\) proved once;
\item monotonicity blocks in which the same floor candidate is controlled by endpoint
inequalities;
\item convexity bounds for \(m^\theta\) over integer intervals;
\item fixed-point integer arithmetic after scaling logarithmic intervals;
\item a hierarchy of block certificates, allowing local rechecking and parallel generation.
\end{enumerate}

A checker should verify the mathematical implication ``certificate \(\Rightarrow\) no
integer value'' with no floating-point operations.  External MPFR code may generate the
certificate but is then removed from the theorem's trust base.  The existing \(4096\)
modules are the prototype; the open engineering problem is certificate compression.

\subsection{Priority 3: exploit the uniform smoothness trichotomy}

\tagProposal Split every hypothetical counterexample into the three cases of
\cref{thm:uniform-second-level}.

\paragraph{Case I: the base-$2$ output is three-smooth.}
Then the primitive odd core is \(w=3\), so
\[
  \beta=a+d\theta,
  \qquad
  M=2^a3^d,
\]
and
\[
  A=3^\beta=3^a(3^\theta)^d\in\N.
\]
Thus \(3^\theta\) is a positive radical whose \(d\)-th power lies in \(\Z[1/3]\).
This is the sharp exceptional target
\[
  3^{\log_2 3}.
\]
The output normalization, exact radical degree, and gcd condition should be specialized to
this one-core case; they may impose stronger congruence or ramification restrictions than
are visible in the generic notation.

\paragraph{Case II: the base-$3$ output is three-smooth.}
Symmetrically the three-free primitive core is \(v=2\), and
\[
  \beta=b+\frac e\theta,
  \qquad
  A=3^b2^e,
\]
so \(2^{1/\theta}=2^{\log_3 2}\) is a positive radical of a rational with denominator
supported at \(2\).  This case deserves a formally symmetric normalization theorem;
currently most of the repository is presented from the base-2 candidate side.

\paragraph{Case III: neither output is three-smooth.}
Then for every pair of noninteger solutions, both \(2^{xy}\) and \(3^{xy}\) are
transcendental.  This is the generic case.  Its advantage is two independent outside-prime
supports; a mixed Archimedean/non-Archimedean determinant may exploit both simultaneously.

The three cases have genuinely different arithmetic.  Treating them separately is more
promising than carrying a generic primitive core through every argument.

\subsection{Priority 4: combine the Archimedean determinant with prime valuations}

\tagProposal The current determinant lower bound uses only that a nonzero integer has
absolute value at least one.  Candidate outputs contain much more information: their
valuations at primes in \(\supp(w)\) and \(\supp(v)\) are exact linear functions of the
coordinate \(k\).  A future determinant should be designed so that column or row
differences force a large predictable divisor while interpolation still makes its
Archimedean absolute value small.

A concrete experimental template is:
\begin{enumerate}
\item choose entries from \(2^{n+k\beta}=2^nM^k\) and
\(3^{n+k\beta}=3^nA^k\), or from monomials in both;
\item arrange finite differences in \(k\) so that powers of \(M\), \(A\), or factors
\(M^h-1\), \(A^h-1\) divide minors;
\item use lifting-the-exponent lemmas at primes from fixed supports to quantify the
non-Archimedean gain;
\item normalize away trivial row factors before estimating the Archimedean determinant;
\item seek a product-formula contradiction rather than the elementary \(|D|\ge1\).
\end{enumerate}

The danger is that the same large factors also enlarge the Archimedean determinant and
cancel the gain.  Small symbolic and numerical models should be used to calculate the exact
balance before investing in formalization.  The target is a genuine improvement over the
three-by-two dimensional count, not a renamed form of the existing determinant.

\subsection{Priority 5: local Kummer analysis of the exact radical}

\tagProposal The canonical radical
\[
  E=w^\theta,
  \qquad E^d=A/3^a,
  \qquad [\Q(E):\Q]=d
\]
produces a pure extension with an exact rational-power index.  Apply local Kummer and
binomial irreducibility theory prime by prime:
\begin{itemize}
\item classify possible ramification at primes dividing \(d\), \(A\), \(w\), and \(3\);
\item combine the power-primitivity of \(w\) with Capelli-type conditions on
\(X^d-A/3^a\);
\item compare the splitting fields of \(E\) and the multiplicative relation
\(E=3^{\log_2w}\);
\item search for a local obstruction in the exceptional cases \(w=3\) or \(v=2\);
\item formalize exact denominator and valuation consequences even when the global
transcendence step remains open.
\end{itemize}

Pure-extension theory alone is unlikely to prove the conjecture, because arbitrary positive
rational radicals certainly exist.  The extra datum is the logarithmic identity
\eqref{eq:bilinear-logs}.  The research question is whether local restrictions on the
radical can be made incompatible with that specific Archimedean logarithmic origin.

\subsection{Priority 6: generalize the determinant from integer to algebraic entries}

\tagProposal Formalize a six-exponentials theorem with algebraic, rather than merely
integral, matrix entries.  If the entries lie in a number field \(K\), a nonzero determinant
\(D\in K\) can be bounded below through
\[
  1\le |N_{K/\Q}(\mathfrak dD)|
\]
after clearing a controlled denominator \(\mathfrak d\).  The conjugates and heights then
enter the upper bound.  This is technically substantial, but it would:
\begin{itemize}
\item formalize \cref{lem:third-base-trans} and all algebraicity equivalences;
\item produce a reusable Mathlib-level six-exponentials component;
\item clarify exactly which height losses prevent shrinking from three logarithms to two;
\item create an interface for mixed algebraic auxiliary values, not only integers.
\end{itemize}

Even if this does not settle Alaoglu--Erd\H{o}s, it would turn the current informal
transcendence upgrade into a kernel-checked theorem and expose the quantitative boundary in
a reusable form.

\subsection{Priority 7: support-first and normalization-first computation}

\tagProposal Raw scanning treats each \(m\) independently.  Under failure, every
noninteger candidate has the form
\[
  m=2^{n+ak}w^{dk}
\]
for one fixed primitive support.  A structured search should enumerate:
\begin{enumerate}
\item a primitive odd core \(w\) (valuation gcd one);
\item the least rational-power index \(d\);
\item the reduced depth \(a\);
\item the nearest integer \(A\) to \(3^a w^{d\theta}\);
\item the simultaneous base-3 primitive normalization and gcd obstruction.
\end{enumerate}

Cheap filters include parity, perfect-power primitivity, denominator support, the exact
common-coordinate gcd condition, and the three-smooth-output trichotomy.  For a bounded
range this is equivalent to the raw search, but it groups inputs by the hypothetical global
structure and may reveal forbidden support patterns.  It also produces data aligned with
the Lean theorems rather than a flat list of near misses.

\subsection{Priority 8: seek arithmetic forcing for finite differences}

\tagProposal The optimized higher-difference theorem is already near its analytic limit.
Progress now requires a mechanism that \emph{forces} many candidates into a short additive
pattern.  Potential sources are:
\begin{itemize}
\item congruences among candidates with the same \(k\)-coordinate;
\item gcd/lcm lattice rectangles whose ordinary integer combinations have controlled size;
\item valuation pigeonhole principles inside a dyadic block;
\item polynomial rather than arithmetic progressions, paired with a matching integer-valued
difference operator;
\item mixed differences in \((n,k)\)-coordinates whose image under
\(m=2^nM^k\) has stronger curvature constants.
\end{itemize}

Every proposal should be tested against the exact conditional monoid.  If the desired
pattern is not forced in a generic free monoid \(\{2^nM^k\}\), then density alone will not
supply it.

\subsection{Priority 9: isolate a nonlinear algebraicity bridge}

\tagProposal By \cref{thm:algebraicity-equivalences}, any proof that
\(6^{\beta^2}\) is algebraic from the original hypotheses completes the conjecture.
Equivalent targets are algebraicity of both \(2^{\beta^2},3^{\beta^2}\), or integrality of
one non-three-smooth base raised to \(\beta\).  Possible sources of such a bridge include:
\begin{itemize}
\item an additional functional equation extracted from leastness or primitive
normalization;
\item an algebraic relation among \(M,A\) forced by matching prime valuations;
\item a determinant whose vanishing yields algebraicity of a second iterate;
\item a rigidity theorem for algebraic values of \(z\mapsto z^\beta\) on a finitely
generated multiplicative group stronger than six exponentials.
\end{itemize}

This is the highest-risk direction but the one most directly aligned with the exact missing
step.  The closure and transcendence theorems in this report provide clean tests: any
claimed bridge must contradict the conditional monoid rather than merely restate it.

\section{Conclusions}\label{sec:conclusion}

The Alaoglu--Erd\H{o}s conjecture survives this investigation.  The existing Lean
formalization already places a hypothetical counterexample in an extraordinarily narrow
shape: a transcendental least generator, a unique primitive odd core, an exact radical
degree, a free rank-two solution monoid, fixed output supports, and strong local
progression exclusions.  None of these constraints alone is inconsistent.

The new results sharpen that picture in directions not present in the starting report.
A counterexample world would have exact rational-function rigidity: the only rational
expressions in the generator that return to the solution set are its natural affine
combinations.  Products of two noninteger solutions never return; quotients return only as
natural multiples; common output powers are measured exactly by coordinate gcd; and
candidate divisibility forms a distributive lattice with a computable ambient saturation.
Six exponentials upgrades the nonlinear failure from nonintegrality to transcendence:
\(6^{xy}\) is transcendental for every pair of noninteger solutions.

These are meaningful advances, but they also make the final boundary cleaner.  To prove the
conjecture one must either:
\begin{enumerate}
\item derive one nonlinear algebraic value from \(2^x,3^x\in\N\), thereby contradicting
the new transcendence dichotomy;
\item manufacture a non-three-smooth integer base with integral \(x\)-th power;
\item find a local valuation obstruction to the exact canonical radical; or
\item build a determinant or product-formula argument that uses integrality beyond the
generic four-exponentials setup.
\end{enumerate}

The finite computation moves the explicit software-level frontier from \(x>20\) to
\(x>26\), while the variable-order analysis explains why finite differences need new
arithmetic forcing.  The most immediate durable gains are to formalize the elementary new
theorems and convert the MPFR scan into replayable certificates.  The most promising
mathematical split is the uniform three-case classification according to which least
outputs are three-smooth, with separate local and transcendence attacks in each case.

No hidden assumption in this report turns these partial results into a proof.  What has
improved is the map of the remaining territory: the conjecture is now surrounded by a
dense set of exact, checkable equivalences and obstructions, and the missing ingredient can
be stated with considerably greater precision than before.

\appendix

\section{Dependency map and Lean-oriented proof sketches}\label{app:lean-sketches}

This appendix records the shortest formalization routes for the new paper-level theorems.
The snippets are schematic and are not claimed to compile unchanged; theorem names from the
repository are shown where known.

\subsection{Dependency map}

\begin{center}
\begin{tabularx}{0.97\textwidth}{@{}lYY@{}}
\toprule
new result & principal existing input & additional mathematics\\
\midrule
rational-function rigidity & canonical primitive generator; transcendence of \(\beta\) & injectivity of polynomial evaluation at a transcendental element\\
product/power criteria & unique \((n,k)\)-coordinates & quadratic coefficient comparison\\
quotient criterion & unique coordinates & denominator clearing and coefficient comparison\\
three-smooth output test & transcendence of \(\theta\) & unique factorization in \(\N\)\\
product transcendence & three-smooth-base theorem & classical six exponentials\\
common power degree & unique coordinates & positive root uniqueness; divisibility\\
primitive statistics & total lattice count & M\"obius inversion\\
gcd/lcm lattice & exact candidate monoid & prime-valuation min/max\\
ambient root saturation & primitive valuation gcd of \(w\) & B\'ezout and linear cone arithmetic\\
optimized differences & arbitrary-order Lean theorem & gamma identity and Stirling\\
\bottomrule
\end{tabularx}
\end{center}

\subsection{Multiplication skeleton}

A direct statement can use the repository predicate
\lean{TwoBaseIntegralSolution}.

\begin{lstlisting}[language={}]
theorem product_solution_iff_integer_factor
    {x y : R}
    (hx : TwoBaseIntegralSolution x)
    (hy : TwoBaseIntegralSolution y) :
    TwoBaseIntegralSolution (x * y) <->
      x in Set.range ((^) : Z -> R) /\
      y in Set.range ((^) : Z -> R) := by
  classical
  constructor
  · intro hxy
    by_cases hAE : AlaogluErdosConjecture
    · left
      exact hAE x hx
    · obtain <w, d, a, c, beta, ..., hcoord, ...> :=
        exists_canonical_primitiveGenerator_of_not_alaogluErdosConjecture hAE
      obtain <nx, kx, hxcoord> := (hcoord x).mp hx
      obtain <ny, ky, hycoord> := (hcoord y).mp hy
      obtain <np, kp, hpcoord> := (hcoord (x * y)).mp hxy
      -- If kx,ky > 0, expansion gives a nonzero quadratic
      -- polynomial over Q vanishing at beta, contradicting
      -- transcendental beta.  Therefore kx = 0 or ky = 0.
      ...
  · rintro (hxint | hyint)
    · exact solution_nsmul_left hxint hy
    · exact solution_nsmul_right hyint hx
\end{lstlisting}

In the actual theorem, the integer range can be simplified to \(\Nzero\), because every
solution is nonnegative.  The key algebraic lemma should be phrased independently:

\begin{lstlisting}[language={}]
lemma transcendental_quadratic_coeff_eq_zero
    {beta a b c : R}
    (hbeta : Transcendental Q beta)
    (h : a * beta^2 + b * beta + c = 0) :
    a = 0 /\ b = 0 /\ c = 0 := ...
\end{lstlisting}

Mathlib's polynomial evaluation API may make it cleaner to construct
\(aX^2+bX+c\) and apply the definition of transcendence.

\subsection{Output smoothness skeleton}

Define a natural predicate
\[
  \operatorname{ThreeSmooth}(m):\Longleftrightarrow
  \exists u,v:\N,\ m=2^u3^v.
\]
Then prove:

\begin{lstlisting}[language={}]
theorem integer_iff_outputs_threeSmooth
    {x : R} (hx : TwoBaseIntegralSolution x) :
    x in Set.range ((^) : Z -> R) <->
      ThreeSmooth (natOutputTwo hx) /\
      ThreeSmooth (natOutputThree hx) := by
  constructor
  · rintro <z, rfl>
    ...
  · rintro <<a,b,h2>, <c,d,h3>>
    have hpoly :
      b * logThreeDivLogTwo^2 +
      (a-d) * logThreeDivLogTwo - c = 0 := by
      -- take logarithms of h2 and h3
      ...
    have htheta := transcendental_logThreeDivLogTwo
    have : b = 0 /\ a-d = 0 /\ c = 0 :=
      transcendental_quadratic_coeff_eq_zero htheta hpoly
    ...
\end{lstlisting}

A practical implementation should avoid partial ``natural output'' definitions by carrying
explicit witnesses \(M,A:\N\) for the two integer powers.

\subsection{Common-power skeleton}

The formal theorem is most reusable if stated for explicit output witnesses.

\begin{lstlisting}[language={}]
theorem simultaneous_perfectPower_iff_divides_coordinates
    (hfail : not AlaogluErdosConjecture)
    {x beta : R} {n k r : N}
    (hr : 0 < r)
    (hbeta : CanonicalPrimitiveGenerator beta)
    (hxcoord : x = n + k * beta) :
    (Exists fun U : N => (U : R)^r = (2 : R)^x) /\
    (Exists fun V : N => (V : R)^r = (3 : R)^x) <->
      r dvd n /\ r dvd k := by
  -- Convert perfect powers to a solution at x/r,
  -- use unique coordinates, and compare.
  ...
\end{lstlisting}

The positive-root step can use strict monotonicity of \(t\mapsto t^r\) on nonnegative
reals or logarithmic injectivity.

\subsection{Ambient saturation skeleton}

The only nontrivial arithmetic lemma is the primitive valuation reconstruction:

\begin{lstlisting}[language={}]
lemma primitive_core_root
    {w u q r : N}
    (hw : oddPrimeValuationGCD w = 1)
    (hr : 0 < r)
    (hpow : u^r = w^q) :
    Exists t : N, q = r*t /\ u = w^t := by
  -- Compare every prime factorization coordinate.
  -- Use a Bezout linear combination of the positive valuations of w
  -- to prove r | q, then use factorization extensionality.
  ...
\end{lstlisting}

The present report's proof uses \(q=dk\).  A general lemma of this form would also be
useful in \lean{PrimitiveOddCore.lean} and \lean{OutputNormalization.lean}.

\subsection{Suggested theorem-order for a pull request}

A low-conflict sequence is:
\begin{enumerate}
\item add a small \lean{ThreeSmooth} API and output-smoothness theorem;
\item add coordinate comparison lemmas for affine, quadratic, and quotient expressions;
\item prove product, power, and quotient criteria;
\item prove common-power and affine-descent criteria;
\item add candidate divisibility and support results;
\item add primitive-core root reconstruction and ambient saturation;
\item only then add asymptotic and six-exponentials extensions.
\end{enumerate}
This ordering keeps the first several commits elementary and independently reviewable.

\section{MPFR verifier source}\label{app:scanner}

The following is the exact source whose SHA-256 digest is given in
\cref{sec:computation}.

\begin{lstlisting}[language=C]
#define _GNU_SOURCE
#include <stdio.h>
#include <stdlib.h>
#include <stdint.h>
#include <inttypes.h>
#include <math.h>
#include <omp.h>

#if defined(__has_include)
#  if __has_include(<mpfr.h>)
#    include <mpfr.h>
#    define HAVE_SYSTEM_MPFR_HEADER 1
#  endif
#endif

#ifndef HAVE_SYSTEM_MPFR_HEADER
/* Minimal declarations for the MPFR 4.x ABI.  The execution environment had
   libmpfr.so.6 but not the development header package.  A normal build with
   libmpfr-dev installed takes the branch above instead. */
typedef unsigned long mp_limb_t;
typedef long mpfr_prec_t;
typedef long mpfr_exp_t;
typedef int mpfr_sign_t;
typedef struct {
  mpfr_prec_t _mpfr_prec;
  mpfr_sign_t _mpfr_sign;
  mpfr_exp_t _mpfr_exp;
  mp_limb_t *_mpfr_d;
} __mpfr_struct;
typedef __mpfr_struct mpfr_t[1];
typedef __mpfr_struct *mpfr_ptr;
typedef const __mpfr_struct *mpfr_srcptr;
typedef enum {
  MPFR_RNDN = 0, MPFR_RNDZ = 1, MPFR_RNDU = 2,
  MPFR_RNDD = 3, MPFR_RNDA = 4, MPFR_RNDF = 5,
  MPFR_RNDNA = -1
} mpfr_rnd_t;
extern void mpfr_init2(mpfr_ptr, mpfr_prec_t);
extern void mpfr_clear(mpfr_ptr);
extern int mpfr_set_ui(mpfr_ptr, unsigned long, mpfr_rnd_t);
extern int mpfr_log(mpfr_ptr, mpfr_srcptr, mpfr_rnd_t);
extern int mpfr_mul(mpfr_ptr, mpfr_srcptr, mpfr_srcptr, mpfr_rnd_t);
extern int mpfr_sub(mpfr_ptr, mpfr_srcptr, mpfr_srcptr, mpfr_rnd_t);
extern int mpfr_cmp_ui(mpfr_srcptr, unsigned long);
extern double mpfr_get_d(mpfr_srcptr, mpfr_rnd_t);
extern const char *mpfr_get_version(void);
#endif

static inline int is_power_of_two_u64(uint64_t x) {
  return x && ((x & (x - 1)) == 0);
}

typedef struct {
  uint64_t checked;
  uint64_t failures;
  double min_lower;
  double min_upper;
  uint64_t min_lower_m, min_lower_n;
  uint64_t min_upper_m, min_upper_n;
  uint64_t first_fail_m, first_fail_guess;
} Stats;

int main(int argc, char **argv) {
  uint64_t start = 1;
  uint64_t limit = (1ULL << 20);
  int bits = 192;
  if (argc > 1) start = strtoull(argv[1], NULL, 0);
  if (argc > 2) limit = strtoull(argv[2], NULL, 0);
  if (argc > 3) bits = atoi(argv[3]);
  if (start < 1 || limit <= start) {
    fprintf(stderr, "bad range\n");
    return 2;
  }

  int nt = omp_get_max_threads();
  Stats *stats = calloc((size_t)nt, sizeof(Stats));
  if (!stats) return 3;
  for (int i = 0; i < nt; ++i) {
    stats[i].min_lower = INFINITY;
    stats[i].min_upper = INFINITY;
  }

  const long double theta = logl(3.0L) / logl(2.0L);
  double t0 = omp_get_wtime();

  mpfr_t two, three, ln2_lo, ln2_hi, ln3_lo, ln3_hi;
  mpfr_init2(two, bits); mpfr_init2(three, bits);
  mpfr_init2(ln2_lo, bits); mpfr_init2(ln2_hi, bits);
  mpfr_init2(ln3_lo, bits); mpfr_init2(ln3_hi, bits);
  mpfr_set_ui(two, 2, MPFR_RNDN);
  mpfr_set_ui(three, 3, MPFR_RNDN);
  mpfr_log(ln2_lo, two, MPFR_RNDD);
  mpfr_log(ln2_hi, two, MPFR_RNDU);
  mpfr_log(ln3_lo, three, MPFR_RNDD);
  mpfr_log(ln3_hi, three, MPFR_RNDU);

#pragma omp parallel
  {
    int tid = omp_get_thread_num();
    Stats *st = &stats[tid];
    mpfr_t xm, xn, xnp1;
    mpfr_t logm_lo, logm_hi, logn_hi, lognp1_lo;
    mpfr_t lhs_lo, lhs_hi, rhs_lo, rhs_hi, dl, du;
    mpfr_init2(xm, bits); mpfr_init2(xn, bits); mpfr_init2(xnp1, bits);
    mpfr_init2(logm_lo, bits); mpfr_init2(logm_hi, bits);
    mpfr_init2(logn_hi, bits); mpfr_init2(lognp1_lo, bits);
    mpfr_init2(lhs_lo, bits); mpfr_init2(lhs_hi, bits);
    mpfr_init2(rhs_lo, bits); mpfr_init2(rhs_hi, bits);
    mpfr_init2(dl, bits); mpfr_init2(du, bits);

#pragma omp for schedule(static)
    for (uint64_t m = start; m < limit; ++m) {
      if (is_power_of_two_u64(m)) continue;

      /* This is only a locator.  The MPFR inequalities below are the proof. */
      uint64_t guess = (uint64_t)floorl(expl(theta * logl((long double)m)));

      mpfr_set_ui(xm, (unsigned long)m, MPFR_RNDN);
      mpfr_log(logm_lo, xm, MPFR_RNDD);
      mpfr_log(logm_hi, xm, MPFR_RNDU);
      mpfr_mul(lhs_lo, logm_lo, ln3_lo, MPFR_RNDD);
      mpfr_mul(lhs_hi, logm_hi, ln3_hi, MPFR_RNDU);

      int ok = 0;
      uint64_t certn = 0;
      double dl_d = 0.0, du_d = 0.0;
      static const int offsets[7] = {0, -1, 1, -2, 2, -3, 3};
      for (int oi = 0; oi < 7; ++oi) {
        int off = offsets[oi];
        if (off < 0 && guess < (uint64_t)(-off) + 1) continue;
        uint64_t n = (uint64_t)((int64_t)guess + off);
        if (n < 1) continue;

        mpfr_set_ui(xn, (unsigned long)n, MPFR_RNDN);
        mpfr_set_ui(xnp1, (unsigned long)(n + 1), MPFR_RNDN);

        /* Directed lower bound for log(m)log(3)-log(n)log(2). */
        mpfr_log(logn_hi, xn, MPFR_RNDU);
        mpfr_mul(rhs_hi, logn_hi, ln2_hi, MPFR_RNDU);
        mpfr_sub(dl, lhs_lo, rhs_hi, MPFR_RNDD);
        if (mpfr_cmp_ui(dl, 0) <= 0) continue;

        /* Directed lower bound for log(n+1)log(2)-log(m)log(3). */
        mpfr_log(lognp1_lo, xnp1, MPFR_RNDD);
        mpfr_mul(rhs_lo, lognp1_lo, ln2_lo, MPFR_RNDD);
        mpfr_sub(du, rhs_lo, lhs_hi, MPFR_RNDD);
        if (mpfr_cmp_ui(du, 0) <= 0) continue;

        ok = 1;
        certn = n;
        dl_d = mpfr_get_d(dl, MPFR_RNDD);
        du_d = mpfr_get_d(du, MPFR_RNDD);
        break;
      }

      if (!ok) {
        ++st->failures;
        if (!st->first_fail_m || m < st->first_fail_m) {
          st->first_fail_m = m;
          st->first_fail_guess = guess;
        }
        continue;
      }

      ++st->checked;
      if (dl_d < st->min_lower) {
        st->min_lower = dl_d;
        st->min_lower_m = m;
        st->min_lower_n = certn;
      }
      if (du_d < st->min_upper) {
        st->min_upper = du_d;
        st->min_upper_m = m;
        st->min_upper_n = certn;
      }
    }

    mpfr_clear(xm); mpfr_clear(xn); mpfr_clear(xnp1);
    mpfr_clear(logm_lo); mpfr_clear(logm_hi);
    mpfr_clear(logn_hi); mpfr_clear(lognp1_lo);
    mpfr_clear(lhs_lo); mpfr_clear(lhs_hi);
    mpfr_clear(rhs_lo); mpfr_clear(rhs_hi);
    mpfr_clear(dl); mpfr_clear(du);
  }

  Stats total = {0};
  total.min_lower = INFINITY;
  total.min_upper = INFINITY;
  for (int i = 0; i < nt; ++i) {
    Stats *s = &stats[i];
    total.checked += s->checked;
    total.failures += s->failures;
    if (s->min_lower < total.min_lower) {
      total.min_lower = s->min_lower;
      total.min_lower_m = s->min_lower_m;
      total.min_lower_n = s->min_lower_n;
    }
    if (s->min_upper < total.min_upper) {
      total.min_upper = s->min_upper;
      total.min_upper_m = s->min_upper_m;
      total.min_upper_n = s->min_upper_n;
    }
    if (s->first_fail_m &&
        (!total.first_fail_m || s->first_fail_m < total.first_fail_m)) {
      total.first_fail_m = s->first_fail_m;
      total.first_fail_guess = s->first_fail_guess;
    }
  }

  double elapsed = omp_get_wtime() - t0;
  printf("mpfr_version=%s\n", mpfr_get_version());
  printf("range=[%" PRIu64 ",%" PRIu64 ")\n", start, limit);
  printf("precision_bits=%d\nthreads=%d\n", bits, nt);
  printf("checked_nonpowers=%" PRIu64 "\n", total.checked);
  printf("failures=%" PRIu64 "\n", total.failures);
  printf("min_lower_log_margin=%.17g at (m,n)=(%" PRIu64 ",%" PRIu64 ")\n",
         total.min_lower, total.min_lower_m, total.min_lower_n);
  printf("min_upper_log_margin=%.17g at (m,n)=(%" PRIu64 ",%" PRIu64 ")\n",
         total.min_upper, total.min_upper_m, total.min_upper_n);
  if (total.failures) {
    printf("first_failure=(m,guess)=(%" PRIu64 ",%" PRIu64 ")\n",
           total.first_fail_m, total.first_fail_guess);
  }
  printf("elapsed_seconds=%.6f\n", elapsed);

  mpfr_clear(two); mpfr_clear(three);
  mpfr_clear(ln2_lo); mpfr_clear(ln2_hi);
  mpfr_clear(ln3_lo); mpfr_clear(ln3_hi);
  free(stats);
  return total.failures ? 1 : 0;
}
\end{lstlisting}

A standard build on a system with MPFR development headers is
\begin{lstlisting}[language=bash]
gcc -O3 -fopenmp mpfr_scan.c -lmpfr -lm -o mpfr_scan
OMP_NUM_THREADS=5 ./mpfr_scan 1 67108864 192
\end{lstlisting}
The recorded audit used sixteen shorter invocations to produce independently inspectable
chunk logs.

\section{Complete chunk audit}\label{app:chunk-audit}

The canonical chunk log has SHA-256 digest
\begin{center}
\small\ttfamily
dcaa7f3b7e5be44acaee9015a2045da14ee42d8fa2d96b9619cdd44c072a6fd7
\end{center}
and contains the full stdout of every invocation.  The table reports the exact scanned
integer interval, number of non-powers certified, and the smallest directed lower and upper
logarithmic margins in each chunk.  Every failure count was zero.

\begin{longtable}{@{}r r r r r@{}}
\caption{Directed-rounding audit through \(2^{26}\).}\label{tab:chunk-audit}\\
\toprule
chunk & integer range & checked & min. lower gap & min. upper gap\\
\midrule
\endfirsthead
\toprule
chunk & integer range & checked & min. lower gap & min. upper gap\\
\midrule
\endhead
0 & 1--4{,}194{,}303 & 4{,}194{,}281 & $1.1539\times10^{-17}$ & $3.0076\times10^{-17}$ \\
1 & 4{,}194{,}304--8{,}388{,}607 & 4{,}194{,}303 & $4.6319\times10^{-20}$ & $1.3515\times10^{-19}$ \\
2 & 8{,}388{,}608--12{,}582{,}911 & 4{,}194{,}303 & $1.2811\times10^{-18}$ & $1.3515\times10^{-19}$ \\
3 & 12{,}582{,}912--16{,}777{,}215 & 4{,}194{,}304 & $4.6319\times10^{-20}$ & $2.9082\times10^{-19}$ \\
4 & 16{,}777{,}216--20{,}971{,}519 & 4{,}194{,}303 & $1.1898\times10^{-18}$ & $1.3515\times10^{-19}$ \\
5 & 20{,}971{,}520--25{,}165{,}823 & 4{,}194{,}304 & $5.4212\times10^{-19}$ & $1.5111\times10^{-19}$ \\
6 & 25{,}165{,}824--29{,}360{,}127 & 4{,}194{,}304 & $4.6319\times10^{-20}$ & $1.7669\times10^{-19}$ \\
7 & 29{,}360{,}128--33{,}554{,}431 & 4{,}194{,}304 & $1.4439\times10^{-19}$ & $5.4132\times10^{-21}$ \\
8 & 33{,}554{,}432--37{,}748{,}735 & 4{,}194{,}303 & $4.2839\times10^{-19}$ & $1.3515\times10^{-19}$ \\
9 & 37{,}748{,}736--41{,}943{,}039 & 4{,}194{,}304 & $3.9292\times10^{-19}$ & $1.0460\times10^{-19}$ \\
10 & 41{,}943{,}040--46{,}137{,}343 & 4{,}194{,}304 & $1.9177\times10^{-19}$ & $1.5111\times10^{-19}$ \\
11 & 46{,}137{,}344--50{,}331{,}647 & 4{,}194{,}304 & $1.0117\times10^{-19}$ & $3.4145\times10^{-19}$ \\
12 & 50{,}331{,}648--54{,}525{,}951 & 4{,}194{,}304 & $5.6927\times10^{-20}$ & $1.2322\times10^{-19}$ \\
13 & 54{,}525{,}952--58{,}720{,}255 & 4{,}194{,}304 & $3.6447\times10^{-20}$ & $1.1345\times10^{-19}$ \\
14 & 58{,}720{,}256--62{,}914{,}559 & 4{,}194{,}304 & $6.7779\times10^{-20}$ & $5.4132\times10^{-21}$ \\
15 & 62{,}914{,}560--67{,}108{,}863 & 4{,}194{,}304 & $5.0923\times10^{-20}$ & $3.9197\times10^{-20}$ \\
\midrule
\multicolumn{2}{r}{\textbf{total}} & \textbf{67{,}108{,}837} & & \\
\bottomrule
\end{longtable}

The global lower-gap record occurred at
\((58{,}570{,}438,2{,}048{,}703{,}434{,}093)\); the global upper-gap record occurred at
\((29{,}411{,}704,687{,}581{,}893{,}443)\).  The repeated upper gap in chunk~14 comes
from exact scaling by two on the input and by three on its \(\theta\)-power.

\section{Reproducibility checklist}\label{app:reproducibility}

To reproduce the finite claim independently:
\begin{enumerate}
\item obtain MPFR~4.2.x and GMP on a 64-bit platform;
\item save the source in \cref{app:scanner} and verify its SHA-256 digest;
\item compile with optimization and OpenMP, or remove OpenMP pragmas for a serial run;
\item run the sixteen half-open intervals listed in \cref{tab:chunk-audit} at 192 bits;
\item verify that every invocation reports \texttt{failures=0};
\item sum \texttt{checked\_nonpowers} and obtain \(67{,}108{,}837\);
\item compare the two global extremal records;
\item optionally rerun at 256 or 320 bits and on a second MPFR build;
\item for a theorem-grade result, translate the intervals into exact rational certificates
and replay them in Lean rather than trusting the executable.
\end{enumerate}

The finite statement is deterministic.  Thread count changes execution time and reduction
order for the recorded minimum metadata, but not whether an individual directed inequality
is certified.

\begin{thebibliography}{99}

\bibitem{AlaogluErdos1944}
L.~Alaoglu and P.~Erd\H{o}s,
\newblock On highly composite and similar numbers,
\newblock \emph{Transactions of the American Mathematical Society} \textbf{56} (1944),
448--469.

\bibitem{Baker1975}
A.~Baker,
\newblock \emph{Transcendental Number Theory},
\newblock Cambridge University Press, 1975.

\bibitem{FousseEtAl2007}
L.~Fousse, G.~Hanrot, V.~Lef\`evre, P.~P\'elissier, and P.~Zimmermann,
\newblock MPFR: A multiple-precision binary floating-point library with correct rounding,
\newblock \emph{ACM Transactions on Mathematical Software} \textbf{33} (2007), no.~2,
article~13.

\bibitem{Gelfond1934}
A.~O.~Gel'fond,
\newblock Sur le septi\`eme probl\`eme de Hilbert,
\newblock \emph{Izvestiya Akademii Nauk SSSR} \textbf{7} (1934), 623--630.

\bibitem{KuipersNiederreiter1974}
L.~Kuipers and H.~Niederreiter,
\newblock \emph{Uniform Distribution of Sequences},
\newblock Wiley, 1974; reprinted by Dover, 2006.

\bibitem{Lang1966}
S.~Lang,
\newblock \emph{Introduction to Transcendental Numbers},
\newblock Addison--Wesley, 1966.

\bibitem{Niven1956}
I.~Niven,
\newblock \emph{Irrational Numbers},
\newblock Carus Mathematical Monographs, Mathematical Association of America, 1956.

\bibitem{Ramachandra1968II}
K.~Ramachandra,
\newblock Contributions to the theory of transcendental numbers, II,
\newblock \emph{Acta Arithmetica} \textbf{14} (1968), 73--88,
\newblock DOI: \url{https://doi.org/10.4064/aa-14-1-73-88}.

\bibitem{Roy1992}
D.~Roy,
\newblock Matrices whose coefficients are linear forms in logarithms,
\newblock \emph{Journal of Number Theory} \textbf{41} (1992), 22--47,
\newblock DOI: \url{https://doi.org/10.1016/0022-314X(92)90081-Y}.

\bibitem{Schneider1934}
T.~Schneider,
\newblock Transzendenzuntersuchungen periodischer Funktionen, I--II,
\newblock \emph{Journal f\"ur die reine und angewandte Mathematik} \textbf{172}
(1934), 65--74 and 175--182.

\bibitem{Waldschmidt2000}
M.~Waldschmidt,
\newblock Open Diophantine problems,
\newblock \emph{Moscow Mathematical Journal} \textbf{4} (2004), 245--305;
preprint \texttt{arXiv:math/0001155} (2000).

\bibitem{TheoremDB2026}
TheoremDB,
\newblock Four exponentials conjecture, problem P4,
\newblock state ``open'', reviewed 31 July 2026,
\newblock \url{https://theoremdb.org/statements/P4/}.

\bibitem{ProveItSnapshot}
V.~Reshetnikov,
\newblock \texttt{ProveIt}: machine-checked mathematics in Lean~4,
\newblock repository snapshot at commit
\commit{5d3982ceb12644fdd3499569d9db5f62f9a31dc6}, 20 August 2026,
\newblock \url{https://github.com/VladimirReshetnikov/ProveIt}.

\bibitem{Mathlib}
The Mathlib Community,
\newblock \emph{Mathlib: the Lean mathematical library},
\newblock \url{https://github.com/leanprover-community/mathlib4}.

\end{thebibliography}

\end{document}
